% info-S1-intro.tex

%-------------------------------------------------------------------------
\section{Contexte scientifique}\label{contexte}
%-------------------------------------------------------------------------
\marginpar{\footnotesize\em
\begin{fig}\label{fig:dico1}[Définitions de l'Académie (1)]
\framebox[8cm]{\begin{minipage}{7.5cm}
{\em INFORMATIQUE} n. f. et adj. XXe siècle. Dérivé d'information sur le modèle de mathématique, électronique.
1. N. f. Science du traitement rationnel et automatique de l'information ; l'ensemble des applications de cette science. 
2. Adj. Qui se rapporte à l'informatique. Système informatique, ensemble des moyens qui permettent de conserver, 
de traiter et de transmettre l'information.\\
\mbox{}\\
{\em INSTRUCTION} n. f. XIVe siècle. Emprunté du latin instructio, «~action d'adapter, de ranger~», 
puis «~instruction~».
Ordre, indication qu'on donne à quelqu'un pour la conduite d'une affaire ; directive, consigne. 
Le plus souvent au pluriel. Des instructions verbales, écrites. Donner des instructions à ses collaborateurs. 
Instructions pour la mise en marche d'un appareil. Par anal. INFORM. Consigne formulée dans un langage de 
programmation, selon un code.\\
\mbox{}\\
{\em LOGICIEL} n. m. XXe siècle. Dérivé de logique.
INFORM. Ensemble structuré de programmes remplissant une fonction déterminée, 
permettant l'accomplissement d'une tâche donnée.\\
\mbox{}\\ 
{\em MATÉRIEL} adj. et n. XIIIe siècle. Emprunté du latin materialis, de même sens.
INFORM. Ensemble des éléments physiques employés pour le traitement des données, par opposition aux 
logiciels.\\
\mbox{}\\
{\em ORDINATEUR} n. m. XVe siècle, au sens de «~celui qui institue quelque chose~» ; XXe siècle, au sens actuel. 
Emprunté du latin ordinator, «~celui qui règle, met en ordre ; ordonnateur~».
Équipement informatique comprenant les organes nécessaires à son fonctionnement autonome, 
qui assure, en exécutant les instructions d'un ensemble structuré de programmes, 
le traitement rapide de données codées sous forme numérique qui peuvent être 
conservées et transmises.
\end{minipage}}
\end{fig}} 

%-------------------------------------------------------------------------
\subsection{Informatique}
%-------------------------------------------------------------------------
Le terme {\sc informatique} est un néologisme proposé en 1962 par
Philippe Dreyfus pour caractériser le traitement automatique de l'information :
il est construit sur la contraction de l'expression  «~information automatique~».
Ce terme a été accepté par l'Académie française en avril 1966, et l'informatique devint alors
officiellement la science du traitement automatique de l'information, où l'information
est considérée comme le support des connaissances humaines et des communications 
dans les domaines techniques, économiques et sociaux (figure \ref{fig:dico1}).
Le mot {\sc informatique} n'a pas vraiment d'équivalent aux Etats-Unis où l'on parle de 
{\em Computing Science} (science du calcul) alors que {\em Informatics} est admis par les Britanniques. 

\begin{defin}\label{def:informatique}
L'informatique est la science du traitement automatique de l'information.
\end{defin}

L'informatique traite de deux aspects complémentaires : les programmes immatériels (logiciel, {\em software})
qui décrivent un traitement à réaliser et les machines (matériel, {\em hardware}) qui exécutent ce traitement.
Le matériel est donc l'ensemble des éléments physiques (microprocesseur, mémoire, écran, clavier, disques durs\ldots)
utilisés pour traiter les données. 
Dans ce contexte, l'ordinateur est un terme générique qui désigne un équipement informatique permettant de traiter des 
informations selon des séquences d'instructions (les programmes) qui constituent le logiciel.
Le terme {\sc ordinateur} a été proposé par le philologue Jacques Perret en avril 1955 
en réponse à une demande d'IBM France qui estimait le mot «~calculateur~» ({\em computer}) 
bien trop restrictif en regard des possibilités de ces machines (voir la proposition de 
Jacques Perret en annexe \ref{annexe:perret} page \pageref{annexe:perret}).

\begin{defin}\label{def:materiel}
Le matériel informatique est un ensemble de dispositifs physiques utilisés pour traiter
automatiquement des informations.
\end{defin}
\begin{defin}\label{def:logiciel}
Le logiciel est un ensemble structuré d'instructions décrivant un traitement d'informations
à faire réaliser par un matériel informatique.
\end{defin}


Un ordinateur n'est rien d'autre qu'une machine effectuant des opéra\-tions
simples sur des séquences de signaux électriques, lesquels sont conditionnés de manière à ne pouvoir
prendre que deux états seulement (par exemple un potentiel électrique maximum ou minimum). Ces
séquences de signaux obéissent à une logique binaire du type «~tout ou rien~» et peuvent donc être
considérés conventionnellement comme des suites de nombres 
ne prenant jamais que les deux valeurs 0 et 1 : 
un ordinateur est donc incapable de traiter autre chose que des 
nom\-bres binaires.
\marginpar{\vspace*{-1.5cm}

\footnotesize\em 
\begin{fig}\label{fig:vonneumann}John Von Neumann (1903-1957).\\ 
\framebox[8cm]{\begin{minipage}{3cm}{\includegraphics[width=2.5cm]{vonneumann.jpg}}
\end{minipage}\hfill
\begin{minipage}{4.5cm}
Mathématicien américain d'ori\-gine hongroise : il a donné son nom 
à l'architecture de von Neumann utilisée dans la quasi 
totalité des ordinateurs modernes.
\end{minipage}}
\end{fig}
\begin{fig}\label{fig:archiVN}Principe de l'architecture de Von Neumann.\\
\framebox[8cm]{\includegraphics[width=6cm]{vonneumann.pdf}}
\end{fig}} 
Toute information d'un autre type doit être
convertie, ou codée, en format binaire. C'est vrai non seulement pour les données que l'on
souhaite traiter (les nombres, les textes, les images, les sons, les vidéos, etc.), mais aussi 
pour les programmes, c'est-à-dire les séquences d'instructions que l'on va fournir à la machine 
pour lui dire ce qu'elle doit faire avec ces données.


L'architecture matérielle d'un ordinateur repose sur le modèle de Von Neumann
(figure \ref{fig:vonneumann}) qui distingue classiquement 4 parties (figure \ref{fig:archiVN}):
\begin{enumerate}
\item L'unité arithmétique et logique, ou unité de traitement, 
	effectue les opérations de base.
\item L'unité de contrôle séquence les opérations.
\item La mémoire contient à la fois les données et le programme qui dira à l'unité de contrôle 
	quels calculs faire sur ces données. La mémoire se divise entre mémoire vive volatile 
	(programmes et données en cours de fonctionnement) et mémoire de masse permanente 
	(programmes et données de base de la machine).
\item Les entrées-sorties sont des dispositifs permettant de communiquer avec le monde extérieur
	(écran, clavier, souris\ldots).
\end{enumerate}
Les 2 premières parties sont souvent rassemblées au sein d'une même structure : le micro-processeur
(la «~puce~»), unité centrale de l'ordinateur.

Dans ce cours, nous ne nous intéresserons qu'aux aspects logiciels de l'informatique.


%-------------------------------------------------------------------------
\subsection{Algorithmique}
%-------------------------------------------------------------------------
\begin{ex}\label{ex:modemploi}
Extrait du mode d'emploi d'un télécopieur concernant l'envoi d'un document.
\marginpar{\footnotesize\em
\begin{td}\label{td:plage1}Dessins sur la plage : exécution (1).\\
On cherche à faire dessiner une figure géométrique sur la plage 
à quelqu'un qui a les yeux bandés.

Quelle figure géométrique dessine-t-on en exécutant la suite d'instructions 
ci-dessous ?
\begin{enumerate}
\item avance de 10 pas,
\item tourne à gauche d'un angle de $120^\circ$,
\item avance de 10 pas,
\item tourne à gauche d'un angle de $120^\circ$, 
\item avance de 10 pas.
\end{enumerate}
\end{td}
}
\begin{enumerate}
\item Insérez le document dans le chargeur automatique.
\item Composez le numéro de fax du destinataire à l'aide du pavé numérique.
\item Enfoncez la touche {\sc envoi} pour lancer l'émission.
\end{enumerate}
\end{ex}
Ce mode d'emploi précise comment envoyer un fax. Il est composé d'une suite 
ordonnée d'instructions (insérez\ldots, composez\ldots, enfoncez\ldots) qui
manipulent des données (document, chargeur automatique, numéro de fax, pavé numérique,
touche {\sc envoi}) pour réaliser la tâche désirée (envoi d'un document).
Chacun a déjà été confronté à de tels documents pour faire fonctionner un 
appareil plus ou moins réticent et donc, consciemment ou non, a déjà 
exécuté un algorithme (ie. exécuter la suite d'instructions dans l'ordre annoncé).
\exo{td:plage1}

\begin{defin}
Un algorithme est une suite ordonnée d'instructions qui indique la démarche 
à suivre pour résoudre une série de problèmes équivalents.
\end{defin}

\begin{ex}\label{ex:chemin}
Extrait d'un dialogue entre un touriste égaré et un autochtone.
\marginpar{\footnotesize\em
\begin{td}\label{td:plage2}Dessins sur la plage : conception (1).\\
Faire dessiner une spirale rectangulaire de 5 côtés, le plus petit côté
faisant 2 pas de long et chaque côté fait un pas de plus que le précédent.
\end{td}
}
\begin{itemize}
\item Pourriez-vous m'indiquer le chemin de la gare, s'il vous plait ?
\item Oui bien sûr : vous allez tout droit jusqu'au prochain carrefour,
	vous prenez à gauche au carrefour et ensuite la troisième à droite,
	et vous verrez la gare juste en face de vous.
\item Merci.
\end{itemize}
\end{ex}

Dans ce dialogue, 
la réponse de l'autochtone est la description d'une suite ordonnée
d'instructions (allez tout droit, prenez à gauche, prenez la troisième à droite) qui
manipulent des données (carrefour, rues) pour réaliser la tâche désirée (aller à la gare).
Ici encore, chacun a déjà été confronté à ce genre de situation et donc, consciemment ou non, 
a déjà construit un algorithme dans sa tête (ie. définir la suite d'instructions pour réaliser 
une tâche).
\marginpar{\footnotesize\em
\begin{fig}\label{fig:dico2}Définitions de l'Académie Française (2).\\
\framebox[8cm]{\begin{minipage}{7.5cm}
{\em ALGORITHME} n. m. XIIIe siècle, augorisme. Altération, sous l'influence du grec 
arithmos, «~nombre~», d'algorisme, qui, par l'espagnol, remonte à l'arabe Al-Khuwarizmi, 
surnom d'un mathématicien.
Méthode de calcul qui indique la démarche à suivre pour résoudre une série de 
problèmes équivalents en appliquant dans un ordre précis une suite finie de règles.\\
\mbox{}\\
{\em DONNÉE} n. f. XIIIe siècle, au sens de «~distribution, aumône~» ; XVIIIe siècle, 
comme terme de mathématiques. Participe passé féminin substantivé de donner au sens de «~indiquer, dire~».
1. Fait ou principe indiscuté, ou considéré comme tel, sur lequel se fonde un raisonnement ; 
constatation servant de base à un examen, une recherche, une découverte. Les données de la science. 
Les données de l'expérience.
{\em MATH.} Souvent au pluriel. Chacune des quantités ou propriétés mentionnées dans l'énoncé d'un problème 
et qui permettent de le résoudre.
{\em INFORM.} Représentation d'une information sous une forme conventionnelle adaptée à son exploitation. 
Le traitement automatique des données. Une banque, une base de données.
\end{minipage}}
\end{fig}}
Mais quand on définit un algorithme, celui-ci ne doit contenir que des instructions 
compréhensibles par celui qui devra l'exécuter (des humains dans les 2 exemples précédents).
\exo{td:plage2}

Dans ce cours, nous devrons apprendre à définir des algorithmes pour qu'ils soient 
compré\-hen\-si\-bles --- et donc exécutables --- par un ordinateur.

\begin{defin}
L'algorithmique est la science des algorithmes.
\end{defin}

L'algorithmique s'intéresse à l'art de construire des algorithmes ainsi
qu'à caractériser leur validité, leur robustesse, leur réutilisabilité, 
leur complexité ou leur efficacité.

\begin{defin}
La validité d'un algorithme est son aptitude à réaliser exactement la tâche
pour laquelle il a été conçu.
\end{defin}
\noindent Si l'on reprend l'exemple \ref{ex:chemin} de l'algorithme de recherche du chemin
de la gare, l'étude de sa validité consiste à s'assurer qu'on arrive effectivement 
à la gare en exécutant scrupuleusement les instructions dans l'ordre annoncé.

\begin{defin}
La robustesse d'un algorithme est son aptitude à se protéger de conditions 
anormales d'utilisation.
\end{defin}
\noindent Dans l'exemple \ref{ex:chemin}, la question de la robustesse de l'algorithme
se pose par exemple si le chemin proposé a été pensé pour un piéton, alors que le 
«~touriste égaré~» est en voiture et que la «~troisième à droite~» est en sens interdit.

\begin{defin}
\mbox{}\marginpar{\footnotesize\em
\begin{td}\label{td:prop}Propriétés d'un algorithme.\\
Reprendre le TD \ref{td:plage1} et illustrer la validité, la robustesse, la réutilisabilité, la
complexité et l'efficacité de l'algorithme proposé pour dessiner sur la plage.
\end{td}}
La réutilisabilité d'un algorithme est son aptitude à être réutilisé pour résoudre
des tâches équivalentes à celle pour laquelle il a été conçu.
\end{defin}
\noindent L'algorithme de recherche du chemin de la gare est-il réutilisable tel quel pour se rendre
à la mairie ? A priori non, sauf si la mairie est juste à côté de la gare.


\marginpar{\footnotesize\em
\begin{fig}\label{fig:pb2source}Du problème au code source.\\
\framebox[8cm]{\includegraphics[width=3cm]{informatique.pdf}}
\end{fig}}
\begin{defin}
La complexité d'un algorithme est le nombre d'instructions élémentaires à exécuter pour
réaliser la tâche pour laquelle il a été conçu.
\end{defin}
\noindent Si le «~touriste égaré~» est un piéton, la complexité de l'algorithme de recherche de chemin
peut se compter en nombre de pas pour arriver à la gare.

\begin{defin}
L'efficacité d'un algorithme est son aptitude à utiliser de manière optimale
les ressources du matériel qui l'exécute.
\end{defin}
\noindent N'existerait-il pas un raccourci piétonnier pour arriver plus vite à la gare ?
\exo{td:prop}

L'algorithmique permet ainsi de passer d'un problème à résoudre à un algorithme
qui décrit la démarche de résolution du problème. La programmation a alors pour rôle de traduire cet
algorithme dans un langage «~compréhensible~» par l'ordinateur afin qu'il
puisse exécuter l'algorithme automatiquement (figure \ref{fig:pb2source}).


%-------------------------------------------------------------------------
\subsection{Programmation}
%-------------------------------------------------------------------------
Un algorithme exprime la structure logique d'un programme
informatique et de ce fait est indépendant du langage de programmation 
utilisé. Par contre, la traduction de l'algorithme dans un langage particulier
dépend du langage choisi et sa mise en \oe uvre dépend également
de la plateforme d'exécution. 

	
\marginpar{\footnotesize\em
\begin{fig}\label{fig:dico3}Définitions de l'Académie Française (3).\\
\framebox[8cm]{\begin{minipage}{7.5cm}
{\em LANGAGE}. n. m. XIIe siècle. Dérivé de langue. Système de signes, de symboles, 
élaboré à partir des langues naturelles et 
constituant un code qui les remplace dans certains cas déterminés 
(on parle aussi de langage symbolique).
Le langage chiffré, secret, conventionnel. Le langage mathématique. 
Langage logique, fondé sur la logique formelle. Langage informatique. 
Langage de programmation.
\end{minipage}}
\end{fig}}
La programmation d'un ordinateur consiste à lui «~expliquer~» en détail 
ce qu'il doit faire, en sachant qu'il ne «~comprend~» pas le langage humain, 
mais qu'il peut seulement effectuer un traitement 
automatique sur des séquences de 0 et de 1.
Un programme n'est rien d'autre qu'une suite d'instructions, encodées en 
respectant de manière très stricte un ensemble de conventions fixées à 
l'avance par un langage informatique (figure \ref{fig:dico3}). 
La machine décode alors
ces instructions en associant à chaque «~mot~»
du langage informatique une action précise.
Le programme que nous écrivons dans le langage informatique
à l'aide d'un éditeur (une sorte de traitement de
texte spécialisé) est appelé programme source (ou code source). 

\marginpar{\footnotesize\em
\begin{fig}\label{fig:dico4}Définitions de l'Académie Française (4).\\
\framebox[8cm]{\begin{minipage}{7.5cm}
{\em BIT} n. m. XXe siècle. Mot anglo-américain, contraction de binary digit, 
«~chiffre binaire~». Chacun des deux chiffres, 0 et 1, de la numération binaire. 
En informatique, le bit est l'unité élémentaire d'information appelée aussi élément binaire.\\
\mbox{}\\
{\em OCTET} n. m. XXe siècle. Dérivé savant du latin octo, «~huit~».
Unité d'information composée de huit bits.
\end{minipage}}
\end{fig}}
Le seul «~langage~» 
que l'ordinateur puisse véritablement «~comprendre~» 
est donc très éloigné de ce que nous utilisons nous-mêmes. 
C'est une 
longue suite de 0 et de 1 (les «~bits~», {\em binary digit}) traités par
groupes de 8 (les «~octets~», {\em byte}), 16, 32, ou même 64
(figure \ref{fig:dico4}). 

\begin{defin}
Un bit est un chiffre binaire (0 ou 1). 
C'est l'unité élémentaire d'information.
\end{defin}

\begin{defin}
Un octet  est une unité d'information composée de 8 bits.
\end{defin}

\begin{ex}\label{ex:megaoctet}Description de clé USB.\\
Sur le descriptif commercial d'une clé USB, on peut lire :
\begin{itemize}
\item Type de produit  : Lecteur flash USB
\item Taille du module : 512 Mo
\item Type d'interface : Hi-Speed USB
\end{itemize}
\end{ex}
\marginpar{\footnotesize\em 
\begin{fig}\label{fig:cleUSB}Clé USB ({Universal Serial Bus}).\\
\framebox[8cm]{\begin{minipage}{4cm}
\includegraphics{usb.jpg}
\end{minipage}\hfill
\begin{minipage}{3.5cm}
Dispositif matériel 
con\-te\-nant une mé\-moi\-re de mas\-se 
(une mémoire flash ou un mini disque dur).
\end{minipage}}
\end{fig}}
\noindent La «~taille du module~» est la capacité de stockage de la clé qui vaut ici 512 Mo (mégaoctets).
Le préfixe «~méga~» est un multiplicateur décimal qui vaut $10^6$ mais s'il est bien adapté au calcul décimal,
il l'est moins au calcul binaire car il ne correspond pas à une puissance entière de 2. 
En effet, la puissance $x$ de 2 telle que $2^x = 10^6$
vaut $\displaystyle x = 6\cdot\frac{\log(10)}{\log(2)} \approx 19.93$. On choisit alors la puissance entière de 2
immédiatement supérieure ($2^{20} = 1\,048\,576$) pour définir le Mo : 
1 Mo = 1\,048\,576 octets.
\exo{td:octets}
\marginpar{\footnotesize\em \begin{td}\label{td:octets}Unités d'information.\\
Combien y a-t-il d'octets dans 1 ko (kilooctet), 
1 Go (gigaoctet), 1 To (téraoctet), 1 Po (pétaoctet), 1 Eo (exaoctet),
1 Zo (zettaoctet) et 1 Yo (yottaoctet) ?
\end{td}}

Pour «~parler~» à un ordinateur, il nous faudra donc utiliser des
systèmes de traduction automatiques, capables de convertir en nombres binaires des suites de
caractères formant des mots-clés (anglais en général) qui seront plus significatifs pour nous.
Le système de traduction proprement dit s'appellera interpréteur ou bien compilateur, suivant la
méthode utilisée pour effectuer la traduction (figure \ref{fig:code2exec}). On appellera langage de
programmation un ensemble de mots-clés (choisis arbitrairement) associé à un ensemble de règles
très précises indiquant comment on peut assembler ces mots pour former des «~phrases~» que
l'interpréteur ou le compilateur puisse traduire en langage machine (binaire).

\begin{description}
\item[Compilation :] La compilation consiste à traduire la totalité du code source en une fois. 
Le compilateur lit toutes les lignes du programme source et produit une nouvelle 
suite de codes appelé programme objet (ou code objet). 
Celui-ci peut désormais être exécuté indépendamment du
compilateur et être conservé tel quel dans un fichier («~fichier exécu\-table~»).
Les langages {\sc Ada}, {\sc C}, {\sc C++} et {\sc Fortran} sont des exemples de langages 
compilés.

\item[Interprétation :] L'interprétation consiste à traduire chaque ligne du programme source en quelques instructions
du langage machine, qui sont ensuite directement exécutées au fur et à mesure
(«~au fil de l'eau~»). Aucun programme objet n'est généré.
L'interpréteur doit être utilisé chaque fois que l'on veut faire fonctionner le programme. 
Les langages {\sc Lisp} et {\sc Prolog} sont des exemples de langages interprétés.
\marginpar{\footnotesize\em 
\begin{fig}\label{fig:code2exec}Du code source à son exécution.\\
\framebox[8cm]{\includegraphics[width=7.8cm]{compilateur.pdf}}
\end{fig}}

L'interprétation est idéale lorsque l'on est en phase d'apprentissage du langage, ou en cours
d'expérimentation sur un projet. Avec cette technique, on peut en effet tester immédiatement toute
modification apportée au programme source, sans passer par une phase de compilation qui demande
toujours du temps.
Mais lorsqu'un projet comporte des fonctionnalités complexes qui doivent s'exécuter
rapidement, la compilation est préférable : un programme compilé fonctionnera
toujours plus vite que son homologue interprété, puisque dans cette technique l'ordinateur
n'a plus à (re)traduire chaque instruction en code binaire avant qu'elle puisse être exécutée.

\item[Semi-compilation :] Certains langages tentent de combiner les deux techniques afin de retirer le meilleur de
chacune. C'est le cas des langages {\sc Python} et {\sc Java} par exemple. 
De tels langages
commencent par compiler le code source pour produire un code intermédiaire, similaire à un langage
machine (mais pour une machine virtuelle), que l'on appelle bytecode, lequel sera ensuite transmis 
à un interpréteur pour l'exécution
finale. Du point de vue de l'ordinateur, le bytecode est très facile à interpréter en langage machine.
Cette interprétation sera donc beaucoup plus rapide que celle d'un code source.

Le fait de disposer en permanence d'un interpréteur permet de tester immédiatement n'importe
quel petit morceau de programme. On pourra donc vérifier le bon fonctionnement de chaque
composant d'une application au fur et à mesure de sa construction.
L'interprétation du bytecode compilé n'est pas aussi rapide que celle d'un véritable code machine,
mais elle est satisfaisante pour de très nombreux programmes.
\end{description}

Dans ce cours, nous choisirons d'exprimer directement les algorithmes dans un 
langage informatique opérationnel : le langage {\sc Python} \cite{swinnen}, sans passer par un
langage algorithmique intermédiaire.
\marginpar{\footnotesize\em {\sc Python} : {\tt www.python.org} \includegraphics[width=4cm]{python-logo.jpg}}

%-------------------------------------------------------------------------
\section{Objectifs du cours}
%-------------------------------------------------------------------------
%La pédagogie désigne les méthodes et pratiques d'enseignement et 
%d'éducation ainsi que toutes les qualités requises pour transmettre 
%un savoir quelconque.
%
%La didactique se différencie de la pédagogie par le rôle central des contenus 
%disciplinaires et par sa dimension épistémologique (la nature des connaissances à enseigner).
%
%L'apprentissage est l'acquisition de nouveaux savoirs ou savoir-faire, 
%c'est-à-dire le processus d'acquisition de connaissances, compétences, 
%attitudes ou valeurs, par l'étude, l'expérience ou l'enseignement.


%-------------------------------------------------------------------------
\subsection{Objectifs thématiques}
%-------------------------------------------------------------------------
L'objectif principal des enseignements d'informatique S1 de l'ENIB
est l'acquisition des notions fondamentales de l'algorithmique.
\marginpar{\footnotesize\em
\begin{td}\label{td:python}Première utilisation de {\sc Python}.\\
Se connecter sur un poste de travail d'une salle informatique.
\begin{enumerate}
\item Lancer {\sc Python}.
\item Utiliser {\sc Python} comme une simple calculette.
\item Quitter {\sc Python}.
\end{enumerate}
Ne pas oublier de se déconnecter du poste de travail.
\end{td}
}
Plus précisément, nous étudierons successivement :
\begin{enumerate}
\item les instructions de base permettant de décrire les algorithmes : affectation, tests, boucles;
\item les procédures et les fonctions qui permettent de structurer et de réutiliser les algorithmes;
	on parlera alors d'encapsulation, de préconditions, de portée des variables, de passage de paramètres,
	d'appels de fonctions, de récursivité et de jeux de tests;
\item les structures de données linéaires : tableaux, listes, piles, files, qui améliorent la
	structuration des données manipulées par les algorithmes. A cette occasion, on évaluera
	la complexité et l'efficacité de certains algorithmes utilisant ces structures linéaires.
\end{enumerate}
Ces différentes notions seront mise en \oe uvre à travers l'utilisation du 
langage {\sc Python}.
\exo{td:python}


%-------------------------------------------------------------------------
\subsection{Objectifs pédagogiques}\label{sub:pedago}
%-------------------------------------------------------------------------
\marginpar{\footnotesize\em
\begin{fig}\label{fig:bloom}Taxonomie de Benjamin Bloom (1913-1999),
psychologue américain spécialisé en pédagogie.

\framebox[8cm]{\begin{tabular}{l@{ : }p{5.5cm}}
Connaître   & définir, distinguer, acquérir, identifier, rappeler, reconnaître\ldots \\
Comprendre  & traduire, illustrer, représenter, dire avec ses mots, distinguer, réécrire, 
		réarranger, expliquer, démontrer\ldots \\
Appliquer   & appliquer, généraliser, relier, choisir, développer, utiliser, employer, 
		transférer, classer, restructurer\ldots \\
Analyser    & distinguer, détecter, classer, reconnaître, catégoriser, déduire, discerner, comparer\ldots \\
Synthétiser & écrire, relater, produire, constituer, transmettre, modifier, créer, 
		proposer, planifier, projeter, spécifier, combiner, classer, formuler\ldots \\
Evaluer     & juger, argumenter, valider, décider, comparer\ldots \\
\end{tabular}}
\end{fig}}
Au cours du semestre S1, nous nous positionnerons principalement sur les 3 premiers 
niveaux de la taxonomie de Bloom qui constitue une référence pédagogique pour la 
classification des 
niveaux d'apprentissage (figure \ref{fig:bloom}) : connaissance, compréhension, application. 
Les 3 derniers niveaux seront plutôt abordés au cours 
du semestre S2 (analyse, synthèse, évaluation).
\begin{enumerate}
\item Connaissance : mémorisation et restitution d'informations dans les mêmes termes.
\item Compréhension : restitution du sens des informations dans d'autres termes.
\item Application : utilisation de règles, principes ou algorithmes pour résoudre un problème, les règles 
	n'étant pas fournies dans l'énoncé.
\item Analyse : identification des parties constituantes d'un tout pour en distinguer les idées.
\item Synthèse : réunion ou combinaison des parties pour former un tout.
\item Evaluation : formulation de jugements qualitatifs ou quantitatifs.
\end{enumerate}

Afin de mieux nous situer par rapport aux différents types de pédagogie
associés, nous «~filerons~» une métaphore musicale inspirée de \cite{charles}.
\begin{description}
\item[Pédagogie par objectifs :] Le solfège est l'étude des principes élémentaires
	de la musique et de sa notation : le musicien «~fait ses gammes~» et chaque exercice
	a un objectif précis pour évaluer l'apprentissage du «~langage musical~».
	Il en va de même pour l'informaticien débutant confronté à l'apprentissage d'un 
	langage algorithmique.
\item[Pédagogie par l'exemple :] L'apprentissage des grands classiques 
	\marginpar{\footnotesize\em
	\begin{rem}On trouvera une liste d'algorithmes classiques sur le site 
	du National Institute of Standards and Technology {\em (NIST : {\tt
	www.nist.gov/dads/})}.
	\end{rem}}
	permet au musicien
	de s'approprier les bases du solfège en les appliquant à ces partitions connues et en les
	(re)jouant lui-même. L'informaticien débutant, en (re)codant lui-même des algorithmes bien
	connus, se constituera ainsi une base de réflexes de programmation en «~imitant~» ces
	algorithmes.
\item[Pédagogie de l'erreur :] Les bogues ({\em bugs}) sont à l'informaticien ce que les fausses notes
	sont aux musiciens : des erreurs. 
	Ces erreurs sont nécessaires dans l'acquisition de la connaissance.
	Un élève a progressé si, après s'être trompé,
	il peut reconnaître qu'il s'est trompé,  dire où et pourquoi il s'est trompé,
	et comment il recommencerait sans produire les mêmes erreurs.
\item[Pédagogie par problèmes :] Connaissant «~ses~» classiques et les bases du solfège, le musicien
	devenu plus autonome peut envisager sereinement la création de ses propres morceaux.
	Le développement d'un projet informatique ambitieux sera «~mis en musique~» au semestre S2.
\end{description}

Dans ce cours, nous adopterons ces différentes stratégies pédagogiques sans oublier
qu'en informatique on apprend toujours mieux en faisant par soi-même.

%-------------------------------------------------------------------------
\subsection{Objectifs comportementaux}
%-------------------------------------------------------------------------
Nous cherchons à développer trois «~qualités~» comportementales chez l'informaticien
débu\-tant : la rigueur, la persévérance et l'autonomie.
	\marginpar{\footnotesize\em
	\begin{td}\label{td:erreur}Erreur de syntaxe en {\sc Python}.\\
	On considère la session {\sc Python} suivante :\\[1mm]
	\mbox{}\hfill\begin{minipage}{7cm}\tt
	>>> x = 3\\
        >>> \ y = x\\
        \mbox{}\ \ File "<stdin>", line 1\\
        \mbox{}\ \ \ \ y = x\\
        \mbox{}\ \ \ \ \char`^\\
        SyntaxError: invalid syntax\\
        >>>
	\end{minipage}\\[1mm]
	De quelle erreur de syntaxe s'agit-il ?
	\end{td}}
\begin{description}
\item[Rigueur :] Un ordinateur est une machine qui exécute vite et bien les
	instructions qu'on lui a «~apprises~». Mais elle ne sait pas interpréter
	autre chose : même mineure, une erreur provoque le dysfonctionnement de la 
	machine.
	\begin{ex}Erreur de syntaxe en langage C.\\
	Pour un «~{\tt ;}~» oublié dans un programme {\sc C}, le code source n'est pas compilable
	et le compilateur nous avertit par ce type de message :\\
	
	\noindent\mbox{}\hspace*{1cm}\begin{minipage}[t]{12cm}\tt
	fichier.c: In function `main':\\
	fichier.c:7: error: syntax error before "printf"
	\end{minipage}
	\end{ex}
	Même si un humain peut transiger sur le «~{\tt ;}~», la machine ne le peut pas : 
	l'ordinateur ne se contente pas de l'«~à peu près~».
	La respect des consignes, la précision et l'exactitude sont donc de 
	rigueur en informatique !
	\exo{td:erreur}
\item[Persévérance :] \mbox{}	\marginpar{\footnotesize\em
	\begin{td}\label{td:plage5}Dessins sur la plage : persévérance.\\
	Finir l'algorithme suivant qui cherche à dessiner un losange sur la plage.
	\begin{enumerate}
	\item avance de 10 pas,
	\item tourne à gauche d'un angle de $60^\circ$,
	\end{enumerate}
	\hspace*{0.9cm}\vdots
	\end{td}}
	Face à l'intransigeance de la machine, le débutant est confronté
	à ses nombreuses erreurs (les siennes, pas celles de la machine !) et sa tendance
	naturelle est de passer à autre chose. 
	Mais le papillonnage (ou {\em zapping})
	est une très mauvaise stratégie en informatique : pour s'exécuter correctement,
	un programme doit être finalisé. L'informatique nécessite d'«~aller au bout des
	choses~».
	\exo{td:plage5}
\item[Autonomie :] Programmer soi-même les algorithmes
	qu'on a définis est sans doute le meilleur moyen pour mieux assimiler
	les principales structures algorithmiques et pour mieux comprendre ses
	erreurs en se confrontant à l'intransigeante impartialité de l'ordinateur,
	vérita\-ble «~juge de paix~» des informaticiens.
	\marginpar{\footnotesize\em
	\begin{td}\label{td:autonomie}Autonomie.\\
	Trouver les définitions du mot autonomie et de son contraire (de son antonyme).
	\end{td}}
	Par ailleurs, l'évolution continuelle et soutenue des langages de programmation et des 
	ordinateurs nécessitent de se tenir à jour régulièrement et seule l'autoformation systématique 
	permet de «~ne pas perdre pied~».
	\exo{td:autonomie}
	
	Pratique personnelle et autoformation constituent ainsi deux piliers de
	l'autonomisation nécessaire de l'apprenti informaticien.
\end{description}


%-------------------------------------------------------------------------
\section{Organisation du cours}
%-------------------------------------------------------------------------

%-------------------------------------------------------------------------
\subsection{Présentiel}
%-------------------------------------------------------------------------
Les enseignements d'informatique S1 de l'ENIB sont dispensés lors de 45h
de séances de cours/TD et de séances de laboratoire.
\marginpar{\footnotesize\em \begin{rem} Un semestre s'étale sur 15 semaines d'enseignements, 
8 semaines consacrées au cours et 7 semaines aux TD. \end{rem}}
\begin{itemize}
\item Les cours/TD ont lieu 1 semaine sur 2 à raison de 3h par semaine, soit
24h de cours/TD sur toute la durée du semestre. Ils se déroulent dans une salle
banalisée.
\item Les séances de laboratoire ont lieu 1 semaine sur 2 en alternance avec les 
cours à raison de 3h par semaine, soit 21h de laboratoire dans le semestre. Elles
se déroulent en salle informatique.
\end{itemize}


%-------------------------------------------------------------------------
\subsection{Documents}
%-------------------------------------------------------------------------
Les principaux documents accompagnant les cours sont de 2 types : 
les supports de cours et les notes de cours.
\begin{description}
\item[Support de cours :] \mbox{}\marginpar{\footnotesize\em 
\begin{fig}\label{fig:transparent}Exemple de transparent de cours.\\
\framebox[8cm]{
%\includeslide[width=7.5cm]{informatique<4>}
}
\end{fig}}
	il s'agit de la copie papier des transparents projetés
	en présentiel (figure \ref{fig:transparent}). 
\item[Notes de cours :] il s'agit de notes qui complètent et
	précisent certains points présentés en cours. Ces notes proposent
	également les exercices de travaux dirigés qui sont étudiés
	en cours et au laboratoire. Le document que vous 
	êtes en train de lire constitue les premières notes de cours
	des enseignements d'informatique S1 de l'ENIB. Il y en aura 3 autres 
	sur les sujets suivants :
	\begin{enumerate}
	\item les instructions de base,
	\item les procédures et les fonctions,
	\item les structures de données linéaires.
	\end{enumerate}
\end{description}

Un site {\sc Web} permet de retrouver ces documents au format {\tt pdf} 
({\em Portable Document Format}) ainsi que d'autres documents tels que
le planning prévisionnel des cours et des contrôles, des exemples corrigés 
d'évaluation, les notes aux différents contrôles ou encore des liens vers 
des sites pertinents pour le cours. 
Un forum de discussions professeurs$\leftrightarrow$élèves
est également ouvert sur ce site.
\exo{td:site}
\marginpar{\footnotesize\em 
\begin{td}\label{td:site}Site {\sc Web} d'Informatique S1.\\
Se connecter sur le site {\sc Web} du cours d'informatique S1 de l'ENIB et
vérifier que ces notes de cours sont bien disponibles sur le site
 au format {\tt pdf}.
\end{td}}

La consultation régulière du site est indispensable pour
se tenir au courant des dernières évolutions du cours : en cas d'ambiguïté,
ce sont les informations de ce site qui feront foi.


%-------------------------------------------------------------------------
\subsection{Evaluations des apprentissages}\label{sub:evaluation}
%-------------------------------------------------------------------------
L'évaluation des connaissances et des compétences acquises par les étudiants
repose sur 4 types de contrôle : les contrôles d'attention, les contrôles de TD, 
les contrôles d'autoformation et les contrôles de compétences.
\begin{description}
\item[Contrôle d'attention :] \mbox{}\marginpar{\footnotesize\em 
\begin{td}\label{td:attention1}Exemple de contrôle d'attention (1).\\
Répondre de mémoire aux questions suivantes (ie. sans rechercher les 
solutions dans les pages précédentes).
\begin{enumerate}
\item Quels sont les 3 principaux thèmes informatiques abordés ?
\item Quelles sont les 4 principales stratégies péda\-go\-giques suivies ?
\item Quelles sont les 3 principales qualités comportementales recherchées ?
\end{enumerate}
\end{td}}
	il s'agit d'un QCM (questionnaire à choix multiples) 
	auquel il faut répondre individuellement
	sans document, en 5' en fin de cours, et dont les questions portent 
	sur des points abordés pendant ce cours. Ce type de contrôle teste directement
	l'acquisition de connaissances au sens du «~connaître~» de la classification de
	Bloom (section \ref{sub:pedago} page \pageref{sub:pedago}).
	Il permet d'évaluer «~à chaud~» la capacité d'attention des étudiants et les incite 
	à une présence attentive afin de bénéficier au maximum des heures de présence 
	aux cours. 	
	\exo{td:attention1}
	
	Des exemples de QCM sont systématiquement proposés dans les notes de cours
	comme par exemple celui du TD \ref{td:qcm} page \pageref{td:qcm} pour cette introduction générale.
\item[Contrôle de TD :]\mbox{}\marginpar{\footnotesize\em 
	\begin{td}\label{td:TD}Exemple de contrôle de TD.\\
	Répondre aux questions du TD \ref{td:attention1} situé dans la marge droite
	de la page \pageref{td:attention1}.
	\end{td}}
	il s'agit ici d'inciter les étudiants à préparer
	activement les séances de laboratoire. En début de chaque séance de 
	laboratoire, chaque binôme doit répondre sans document en 10' aux questions
	d'un exercice de TD.
	L'exercice est choisi aléatoirement parmi les exercices de TD situés en marge
	des notes de cours et se rapportant au thème du TD. 
	\exo{td:TD}

	Il concerne le «~comprendre~» et l'«~appliquer~» de la taxonomie de Bloom
	(section \ref{sub:pedago} page \pageref{sub:pedago}).
\item[Contrôle d'autoformation :] \mbox{}\marginpar{\footnotesize\em
	\begin{td}\label{td:bool}Exemple de contrôle d'autoformation (1).\\
	Etablir la table de vérité du circuit logique ci-dessous où $a$, $b$
	et $c$ sont les entrées, $s$ et $t$ les sorties.\\[2mm]
	\centerline{\includegraphics[width=6cm]{add3.pdf}}
	\end{td}}
	un contrôle d'autoformation porte sur un thème
	prévu à l'avance et que l'étudiant doit approfondir par lui-même. Les thèmes 
	retenus pour l'autoformation en S1 sont par exemple le calcul booléen,
	le codage des nombres ou encore la recherche d'éléments dans un tableau de données.
	\exo{td:bool}

	Ces contrôles se déroulent pendant les séances de cours : 
	ce sont des écrits individuels de 30' sans document qui concernent
	le «~connaître~», le «~comprendre~» et l'«~appliquer~» de la taxonomie de Bloom.
\item[Contrôle de compétences :] \mbox{}\marginpar{\footnotesize\em
	\begin{td}\label{td:competences}Exemple de contrôle des compétences.\\
	Répondre aux questions du TD \ref{td:shadok} de la page \pageref{td:shadok}.
	\end{td}}
	les contrôles de compétences (ou DS) durent 80'
	pendant une séance de cours. Plus longs, ils permettent d'aborder l'un des 3 
	derniers niveaux de la classification de Bloom (analyser, synthétiser, évaluer)
	comme le font les exercices de la section \ref{sub:analyser} page \pageref{sub:analyser}.
	\exo{td:competences}
\end{description}

\noindent Quel que soit le type de contrôle, un exercice cherche à évaluer un objectif
particulier. Aussi, la notation exprimera la distance qui reste à parcourir 
pour atteindre cet objectif (figure \ref{fig:cible}): 
\marginpar{\footnotesize\em
\begin{fig}\label{fig:cible}Notation : la métaphore de la cible.\\
\framebox[8cm]{\includegraphics{cible.pdf}}
\end{fig}}
$$\begin{tabular}{l@{ : }l@{ $\rightarrow$ }l}
0 & «~en plein dans le mille !~» & l'objectif est atteint \\
1 & «~pas mal !~» & on se rapproche de l'objectif \\
2 & «~juste au bord de la cible !~» & on est encore loin de l'objectif\\
3 & «~la cible n'est pas touchée !~» & l'objectif n'est pas atteint
\end{tabular}$$
Ainsi, et pour changer de point de vue sur la notation, le contrôle 
est réussi lorsqu'on a 0 ! Il n'y a pas non plus de $1/2$ point ou de $1/4$ 
de point : le seul barême possible ne comporte que 4 niveaux : 0, 1, 2 et 3.
On ne cherche donc pas à «~grappiller~» des points : 
\begin{itemize}
\item on peut avoir 0 (objectif atteint) et avoir fait une ou deux erreurs 
	bénignes en regard de l'objectif recherché;
\item on peut avoir 3 (objectif non atteint) et avoir quelques éléments de
	réponse corrects mais sans grand rapport avec l'objectif.
\end{itemize}


%-------------------------------------------------------------------------
\subsection{Evaluation des enseignements}
%-------------------------------------------------------------------------
En fin de semestre, les étudiants organisent de manière anonyme et 
dans le respect des personnes, une évaluation individuelle des enseignements. 
Elle est structurée en 2 parties :
un questionnaire de 10 à 20 questions (voir un exemple en annexe
\ref{annexe:eval} page \pageref{annexe:eval}) 
auxquelles il faut répondre selon une grille pré-définie
et une partie « expression libre » que l'étudiant remplit, ou non, à son gré.

L'ensemble des fiches d'évaluation est remis à l'équipe pédagogique
d'Informatique S1, qui après en avoir pris connaissance, organise une rencontre spécifique
avec les étudiants. Cette évaluation a pour objectif l'amélioration de la qualité 
des enseignements et de la pédagogie à partir de la perception qu'en ont les étudiants.


%-------------------------------------------------------------------------
\section{Méthodes de travail}
%-------------------------------------------------------------------------
Il ne s'agit pas ici d'imposer des méthodes de travail, 
mais de fournir des pistes pour ceux qui en cherchent.
Dans tous les cas, l'expérience montre que :
\begin{enumerate}
\item la seule présence, même attentive et active, aux séances de cours 
	et de laboratoire ne suffit pas : il faut prévoir un temps de travail personnel
	qui, selon l'étudiant et la matière, peut aller de 50\% à 150\% du temps de 
	présence en cours;
\item la régularité dans le travail personnel est un gage d'apprentissage
	plus efficace.
\end{enumerate}
\marginpar{\footnotesize\em
\begin{td}\label{td:controles}Nombre de contrôles.\\
Après consultation du calendrier prévisionnel de l'annexe \ref{annexe:planning} 
page \pageref{annexe:planning}, donner le nombre et le type de contrôles prévus
au calendrier du semestre.
\end{td}}
Le calendrier prévisionnel des enseignements et des contrôles associés
est distribué en début de semestre (un exemple de planning est donné
en annexe \ref{annexe:planning} page \pageref{annexe:planning}).
\exo{td:controles} 

\noindent Les documents de cours sont distribués au moins 15 jours avant
les séances correspondantes (sauf en ce qui concerne les 2 premières semaines
de cours évidemment).
Par ailleurs, à tout moment, le calendrier et les documents sont disponibles 
sur le site {\sc Web} d'Informatique S1.

%-------------------------------------------------------------------------
\subsection{Avant, pendant et après le cours}\label{sub:cours}
%-------------------------------------------------------------------------
\begin{description}
\item[Préparation :] Certains cours débutent par un contrôle d'autoformation
	(voir section \ref{sub:evaluation}) dont le thème est en rapport avec certains
	aspects du cours qui suivra immédiatement.
	Il est donc nécessaire d'avoir étudié avant et par soi-même le
	sujet du contrôle.
	
	En général, on trouvera les informations nécessaires
	soit directement sur le site d'Informatique S1, 
	soit dans les références bibliographiques données en fin des notes de cours
	(voir par exemple les références de ce chapitre en page \pageref{biblio}),
	soit dans les polycopiés d'autres cours de l'ENIB (mathématiques, 
	électronique, productique, microprocesseurs\ldots), soit encore sur internet
	en s'assurant de la qualité du site (préférer des sites universitaires ou 
	d'écoles dont l'activité principale est d'enseigner ces matières).
	
	\marginpar{\footnotesize\em
	\begin{td}\label{td:negation}Exemple de contrôle d'autoformation (2).\\
	Exprimer en la développant la négation des expressions booléennes suivantes :
	\begin{enumerate}
	\item {\tt (0 < x) and (x < 3)}
	\item {\tt (x < -2) or (x > 4)}
	\item {\tt a and (not b)}
	\item {\tt (not a) or b}
	\end{enumerate}
	\end{td}}
	Par exemple, il est nécessaire d'avoir assimilé les principes
	du calcul booléen pour maîtriser les tests et les boucles du
	langage algorithmique. C'est pourquoi, une autoformation est
	imposée sur ce domaine déjà connu des étudiants.
	Un contrôle-type d'autoformation sur le calcul booléen 
	pourrait par exemple être composé des TD \ref{td:bool} 
	page \pageref{td:bool} et \ref{td:negation} ci-contre.\\
	\exo{td:negation}
	
	Pour chaque contrôle d'autoformation, un exemple corrigé est
	disponible sur le site d'Informatique S1. Il est donc fortement 
	recommandé de travailler ce contrôle-type : après avoir revu
	par soi-même les principales notions du domaine, faire
	le contrôle sans consulter au préalable la solution proposée. 
	\marginpar{\footnotesize\em
	\begin{td}\label{td:attention2}Exemple de contrôle d'attention (2).\\
	Répondre de mémoire aux questions suivantes (ie. sans rechercher les 
	solutions dans les pages précédentes).
	\begin{enumerate}
	\item Quels sont les 4 types de contrôle proposés ?
	\item Quels sont les documents que l'on peut trouver sur le site {\sc Web}
		du cours ?
	\end{enumerate}
	\end{td}}
\item[Participation :] Par respect pour les autres participants, 
	la ponctualité est de rigueur pour l'étudiant comme pour le professeur.
	Mais assister à un cours n'a d'intérêt que dans le cas d'une 
	présence attentive et soutenue : le temps passé en cours 
	pour assimiler les nouvelles notions sera gagné sur le temps de
	travail personnel futur nécessaire à leur assimilation.
	Chaque page des supports de cours est constituée
	d'une copie d'un transparent de cours dans la demi-page supérieure
	alors que la demi-page inférieure est volontairement laissée vierge
	pour que l'étudiant prenne des notes pendant le cours (figure \ref{fig:support}).
	\marginpar{\footnotesize\em
	\begin{fig}\label{fig:support}Exemple de support de cours.\\
	\framebox[8cm]{
	\framebox[4.5cm]{\begin{minipage}{4cm}
	%\includeslide[width=3.8cm]{informatique<4>}
	\vspace*{3cm}
	
	\end{minipage}}}
	\end{fig}}
	La prise de note systématique est en effet un facteur qui favorise l'attention.
	Un contrôle de type QCM en fin de cours évaluera l'attention
	des étudiants (voir section \ref{sub:evaluation}). 
	\exo{td:attention2}

	Le cours est illustré par des exemples et des exercices : à ces occasions, 
	c'est une participation active de l'étudiant qui est attendue dans une
	démarche volontaire d'assimilation du cours « à chaud ».
	
\item[Appropriation :] Dans la mesure du possible, il faut relire ses propres notes 
	ainsi que les notes de cours le soir même afin de « fixer » les principales
	idées du cours. La révision proprement dite peut venir ultérieurement
	quelques jours plus tard (et non quelques semaines ni quelques mois après).
	
	Les définitions, comme toute définition, sont à apprendre par c\oe ur.
	Une technique possible est de lire 2 fois de suite à voix haute la définition
	en détachant distinctement les différentes expressions (exemple: «~l'informatique~»
	«~est la science~» «~du traitement~» «~automatique~» «~de l'information~»),
	puis l'écrire de mémoire.
	
	Pour réviser le cours, faire systématiquement les TD au moment où ils sont
	référencés dans les notes par le symbole $\rule{2mm}{2mm}\,$. 
	C'est particulièrement vrai 
	avec les exemples pour lesquels sont associés des exercices en marge des
	notes de cours
	(exemple : l'exemple \ref{ex:modemploi} de la page \pageref{ex:modemploi}
	et le TD \ref {td:plage1} en marge droite de la même page).
	Lorsque l'exemple est un algorithme, il faut  systématiquement se mettre 
	mentalement à la place de la machine qui va les exécuter (on parle alors 
	d'«~empathie numérique~») afin de vérifier le résultat obtenu. 
	Pour cela, il faut être méthodique et rigoureux. 
	Et petit à petit, à force de pratique, l'expérience fera qu'on « verra » 
	le résultat produit par les instructions au fur et à mesure qu'on les écrit. 
	Naturellement, cet apprentissage est long, et demande des heures de travail 
	patient. Aussi, dans un premier temps, il faut éviter de sauter les étapes : 
	la vérification méthodique, pas à pas, de chacun des algorithmes 
	représente plus de la moitié du travail à accomplir \cite{darmengeat}.
\end{description}

%-------------------------------------------------------------------------
\subsection{Avant, pendant et après le laboratoire}
%-------------------------------------------------------------------------
\begin{description}
\item[Préparation :] \mbox{}	 \marginpar{\footnotesize\em
	 \begin{td}\label{td:exercices}Nombres d'exercices de TD.\\
	 Combien d'exercices y avait-il à faire avant celui-ci ?
	 \end{td}}
	Faire les exercices situés dans les marges de ces notes de cours
	est non seulement une façon de réviser son cours (voir section \ref{sub:cours})
	mais aussi de préparer les séances de laboratoire.
	Pour ces notes de cours, la liste complète des exercices est donnée en 
	page \pageref{sub:listexo}.
	Un de ces exercices choisi aléatoirement fera l'objet d'une
	évaluation en début de chaque séance de TD (voir section \ref{sub:evaluation}).
	\exo{td:exercices}
	
	Tous ces exercices ne nécessitent pas une longue phase de réflexion comme
	les TD \ref{td:prop} ou \ref{td:competences} ($\rightarrow$ \ref{td:shadok}).
	Certains au contraire ne présentent aucune difficulté particulière
	comme les TD \ref{td:controles} et \ref{td:exercices}
	qui demandent simplement de compter les contrôles et les exercices.
	D'autres comme les TD \ref{td:attention1} et \ref{td:attention2} 
	font appel à la mémoire, d'autres encore à la recherche d'informations
	comme les TD \ref{td:octets} et \ref{td:autonomie}, ou à une mise en
	\oe uvre pratique comme les TD \ref{td:python} et \ref{td:site}.
\item[Participation :] \mbox{}\marginpar{\footnotesize\em
	\begin{td}\label{td:labo}Environnement de travail.\\
	Sur un poste de travail d'une salle informatique :
	\begin{enumerate}
	\item Quel est le type de clavier ?
	\item Comment ouvre-t-on un terminal ?
	\item Comment lance-t-on {\sc Python} ?
	\item Où sont stockés les fichiers de travail ?
	\end{enumerate}
	\end{td}}
	Ici encore, par respect pour les autres participants, 
	la ponctualité est de rigueur pour l'étudiant comme pour le professeur.

	En informatique, lorsqu'on travaille en binôme, il faut régulièrement 
	alterner les rôles entre l'«~écrivain~» qui manipule clavier et souris 
	et le «~lecteur~» qui vérifie la production de l'écrivain. A la fin
	du semestre, chaque étudiant doit être devenu un «~lecteur-écrivain~».
	La pratique est en effet nécessaire pour «~apprivoiser~» l'environnement 
	de travail afin que la machine devienne «~transparente~» et que seul 
	subsiste le problème algorithmique à résoudre.
	
	Pour certains exercices, la solution est donnée dans les notes de cours
	(en section \ref{sub:solutions} page \pageref{sub:solutions} ici).
	Pendant les séances de laboratoire, il faut donc savoir «~jouer le jeu~»
	pour ne pas simplement «~recopier~» la solution proposée (il existe
	d'ailleurs d'autres solutions pour la plupart des exercices).
	
\marginpar{\footnotesize\em
\begin{fig}\label{fig:abelson}«~L'ivresse joyeuse de la programmation~».\\
\framebox[8cm]{\begin{minipage}{7.5cm}
« Chaque programme d'ordinateur est un modèle, forgé par l'esprit, d'un
processus réel ou imaginaire. Ces processus, qui naissent de l'expérience et
de la pensée de l'homme, sont innombrables et complexes dans leurs détails. 
A tout moment, ils ne peuvent être compris que partiellement. 
Ils ne sont que rarement modélisés d'une façon satisfaisante dans
nos programmes informatiques. Bien que nos programmes soient des ensembles
de symboles ciselés avec soin, des mosaïques de fonctions
entrecroisées, ils ne cessent d'évoluer. 
Nous les modifions au fur et à mesure que notre perception du modèle
s'approfondit, s'étend et se généralise, jusqu'à atteindre un 
équilibre métastable aux frontières d'une autre modélisation
possible du problème. L'ivresse joyeuse qui accompagne la
programmation des ordinateurs provient des allers-retours continuels,
entre l'esprit humain et l'ordinateur, des mécanismes exprimés par
des programmes et de l'explosion de visions nouvelles qu'ils apportent.
Si l'art traduit nos rêves, l'ordinateur les réalise sous la forme de 
programmes\,! »\\\mbox{}\hfill
{\em Abelson H., Sussman G.J. et Sussman J., \cite{abelson}}
\end{minipage}}
\end{fig}}
	Un apprenti programmeur est toujours confronté à ses propres erreurs
	(revoir le TD \ref{td:erreur} page \pageref{td:erreur} par exemple). 
	En général, ce sont des erreurs simples à corriger à condition
	de lire les messages d'erreur et de faire l'effort 
	de les comprendre avant d'appeler le professeur «~à l'aide~».
	\begin{ex}Erreur de nom en {\sc Python}.\\
	Voici une erreur classique d'une variable non correctement initialisée :\\[1mm]
	\mbox{}\hspace*{1cm}\begin{minipage}{12cm}\tt
	>>> print x\\
        Traceback (most recent call last):\\
        \mbox{}\ \ File "<stdin>", line 1, in ?\\
        NameError: name 'x' is not defined\\
        >>> 
	\end{minipage}
	\end{ex}
	En {\sc Python}, la dernière ligne du message d'erreur est la plus 
	«~parlante~» au débutant.
\item[Appropriation :] Avant le cours suivant, il faut refaire les TD qui
	ont posé des problèmes et faire les exercices qui n'ont pas pu
	être abordés en séance de laboratoire (les solutions de ces
	exercices complémentaires sont données dans les notes de cours).
\end{description}	

%-------------------------------------------------------------------------
\newpage
\section{Exercices complémentaires}
%-------------------------------------------------------------------------

%-------------------------------------------------------------------------
\subsection{Connaître}
%-------------------------------------------------------------------------
\begin{td}\label{td:qcm}\em QCM (un seul item correct par question).
\begin{enumerate}
\item L'informatique est la science
\marginpar{\footnotesize\em
\begin{rem} Parmi les 5 items de la question ci-contre, un seul item définit
	l'informatique, les 4 autres définissent d'autres sciences.
	Lesquelles ?
\end{rem}
}
	\begin{enumerate}
	\item des dispositifs dont le fonctionnement dépend de 
		la circulation d'électrons
	\item des signaux électriques porteurs d'information ou d'énergie
	\item du mouvement des systèmes matériels et de leurs déformations
	\item du traitement automatique de l'information
	\item de la commande des appareils fonctionnant sans intervention humaine
	\end{enumerate}
\item Le logiciel est
	\begin{enumerate}
	\item la mémoire de l'ordinateur
	\item l'ensemble de dispositifs physiques utilisés pour traiter
		automatiquement des informations
	\item le traitement automatique de l'information
	\item l'ensemble des données manipulées par les instructions
	\item un ensemble structuré d'instructions décrivant un traitement d'informations
		à faire réaliser par un matériel informatique
	\end{enumerate}
\item L'algorithmique est la science
	\begin{enumerate}
	\item du traitement automatique de l'information
	\item des algorithmes
	\item des langages de programmation
	\item des instructions
	\item des données traitées par les algorithmes
	\end{enumerate}
\item Un algorithme est
	\begin{enumerate}
	\item un ensemble de programmes remplissant une fonction déterminée,
		permettant l'accomplissement d'une tâche donnée
	\item une suite ordonnée d'instructions qui indique la démarche 
		à suivre pour résoudre une série de problèmes équivalents
	\item le nombre d'instructions élémentaires à exécuter pour
		réaliser une tâche donnée
	\item une machine effectuant des opérations simples sur des séquences 
		de signaux électriques
	\item un ensemble de dispositifs physiques utilisés pour traiter
		automatiquement des informations
	\end{enumerate}
\item La validité d'un algorithme est son aptitude
\marginpar{\footnotesize\em
\begin{rem} Parmi les 5 items de la question ci-contre, un seul item définit
	la validité d'un algorithme, les 4 autres se rapportent à d'autres propriétés
	des algorithmes. Lesquelles ?
\end{rem}
}
	\begin{enumerate}
	\item à utiliser de manière optimale les ressources du matériel qui l'exécute
	\item à être réutiliser pour résoudre des tâches équivalentes à celle pour 
		laquelle il a été conçu
	\item à se protéger de conditions anormales d'utilisation
	\item à calculer le nombre d'instructions élémentaires nécessaires pour
		réaliser la tâche pour laquelle il a été conçu
	\item à réaliser exactement la tâche pour laquelle il a été conçu
	\end{enumerate}
\item La complexité d'un algorithme est
	\begin{enumerate}
	\item le nombre de fois où l'algorithme est utilisé dans un programme
	\item le nombre de données manipulées par les instructions de
		l'algorithme
	\item le nombre d'octets occupés en mémoire par l'algorithme
	\item le nombre de dispositifs physiques utilisés pour traiter
		automatiquement les instructions
	\item le nombre d'instructions élémentaires à exécuter pour
		réaliser la tâche pour laquelle il a été conçu
	\end{enumerate}
\item Un bit est 
	\begin{enumerate}
	\item un chiffre binaire
	\item composé de 8 chiffres binaires
	\item un chiffre héxadécimal
	\item un caractère
	\item un mot d'un langage informatique
	\end{enumerate}
\item Un compilateur 
	\begin{enumerate}
	\item exécute le code source
	\item traduit le code objet en bytecode 
	\item exécute le bytecode
	\item traduit un code source en code objet
	\item exécute le code objet
	\end{enumerate}
\end{enumerate}
\end{td}

%-------------------------------------------------------------------------
\subsection{Comprendre}
%-------------------------------------------------------------------------
\begin{td}\label{td:plage3}\em Dessins sur la plage : exécution (2).
\begin{enumerate}
\item Quelle figure géométrique dessine-t-on en exécutant la suite d'instructions 
ci-dessous ?
	\begin{enumerate}
	\item avance de 3 pas,
	\item tourne à gauche d'un angle de $90^\circ$,
	\item avance de 4 pas,
	\item rejoindre le point de départ.
	\end{enumerate}
\item Combien de pas a-t-on fait au total pour rejoindre le point de départ ?
\end{enumerate}
\end{td}

\begin{td}\label{td:plage4}\em Dessins sur la plage : conception (2).\\
Reprendre le TD \ref{td:plage2} et illustrer la validité, la robustesse, 
la réutilisabilité, la complexité et l'efficacité de
l'algorithme proposé pour dessiner une spirale rectangulaire.
\end{td}

\begin{td}\label{td:mips}\em Puissance de calcul.\\
Donner l'ordre de grandeur en instructions par seconde des machines suivantes 
(voir \cite{delahaye} par exemple) :
\begin{enumerate}
\item le premier micro-ordinateur de type PC,
\item une console de jeu actuelle,
\item un micro-ordinateur actuel,
\item {\em Deeper-Blue} : l'ordinateur qui a « battu » Kasparov aux échecs en 1997,
\item le plus puissant ordinateur actuel.
\end{enumerate}
\end{td}

\begin{td}\label{td:stock}\em Stockage de données.\\
Donner l'ordre de grandeur en octets pour stocker en mémoire
(voir \cite{delahaye} par exemple) :
\begin{enumerate}
\item une page d'un livre,
\item une encyclopédie en 20 volumes,
\item une photo couleur,
\item une heure de vidéo,
\item une minute de son,
\item une heure de son.
\end{enumerate}
\end{td}
%-------------------------------------------------------------------------
\subsection{Appliquer}
%-------------------------------------------------------------------------
\marginpar{\footnotesize\em
\begin{fig}\label{fig:polygone}Pentagone, hexagone, octogone.\\
\framebox[8cm]{
\includegraphics[width=2.5cm]{pentagone.jpg}
\hfill
\includegraphics[width=2.5cm]{hexagone.jpg}
\hfill
\includegraphics[width=2.5cm]{octogone.jpg}
}
\end{fig}}
\begin{td}\label{td:tortue}\em Tracés de polygones réguliers.\\
On cherche à faire dessiner une figure polygonale (figure \ref{fig:polygone}) 
sur la plage à quelqu'un qui a les yeux bandés.
Pour cela, on ne dispose que de 2 commandes orales :
avancer de $n$ pas en avant ($n$ est un nombre entier de pas) et 
tourner à gauche d'un angle $\theta$ (rotation sur place de $\theta$).
\begin{enumerate}
\item Faire dessiner un pentagone régulier de 10 pas de côté.
\item Faire dessiner un hexagone régulier de 10 pas de côté.
\item Faire dessiner un octogone régulier de 10 pas de côté.
\item Faire dessiner un polygone régulier de $n$ côtés de 10 pas chacun.
\end{enumerate}
\end{td}

%-------------------------------------------------------------------------
\subsection{Analyser}\label{sub:analyser}
%-------------------------------------------------------------------------
\marginpar{\footnotesize\em
\begin{rem}La multiplication en Egypte antique.\\[1mm]
\begin{minipage}{4cm}
La multiplication «~à la rus\-se~»
est une variante connue d'une technique égyptienne
décrite dans le papyrus Rhind (environ -1650). Le scribe Ahmès
y expose des problèmes de géométrie et d'arithmétique (qui viennent en partie 
des Babyloniens) dont cette technique de multiplication.
\end{minipage}
\hfill
\begin{minipage}{3.6cm}
\includegraphics[width=3.5cm]{papyrus.jpg}
\end{minipage}
\end{rem}}
\begin{td}\label{td:russe}\em La multiplication «~à la russe~».\\
La technique de multiplication dite «~à la rus\-se~» consiste à diviser par 
2 le multiplicateur (et ensuite les quotients obtenus), 
jusqu'à un quotient nul, à noter les restes,
et à multiplier parallèlement le multiplicande par 2. 
On additionne alors les multiples obtenus du multiplicande 
correspondant aux restes non nuls.

\noindent Exemple : \begin{minipage}[t]{10cm}
$68 \times 123\ (= 8364)$\\
\begin{tabular}{|r|r|r|r|}
\hline
multiplicande & multiplicateur & reste      & somme partielle\\
$M \times 2$  & $m \div 2$     & $m \bmod 2$ &       \\
\hline
123  & 68 &   0   &   $(0\times 123) + 0$ \\
246  & 34 &   0   &   $(0\times 246) + 0$ \\
492  & 17 &   1   &   $(1\times 492) + 0$ \\
984  &  8 &   0   &   $(0\times 984) + 492$  \\
1968 &  4 &   0   &   $(0\times 1968) + 492$ \\
3936 &  2 &   0   &   $(0\times 3936) + 492$ \\
7872 &  1 &   1   &   $(1\times 7872) + 492$ \\
\hline
\multicolumn{3}{|r|}{$68 \times 123 =$} & 8364\\
\hline
\end{tabular}
\end{minipage}
\end{td}

\noindent Effectuer les multiplications suivantes selon la technique «~à la russe~».
\begin{enumerate}
\item $64 \times 96\ (= 6144)$
\item $45 \times 239\ (= 10755)$
\end{enumerate}

\begin{td}\label{td:ibnalbanna}\em La multiplication arabe (d'après \cite{chabert}).
\marginpar{\footnotesize\em 
\begin{fig}\label{fig:ibnalbanna}Principe du tableau d'Ibn al-Banna.\\
\framebox[8cm]{\includegraphics[width=7.5cm]{ibnalbanna.pdf}}
\end{fig}}\\
On consid\`ere ici le texte d'Ibn al-Banna concernant la multiplication
\`a l'aide de tableaux.

{\footnotesize
« Tu construis un quadrilat\`ere que tu subdivises verticalement et
horizontalement en autant de bandes qu'il y a de positions dans les
deux nombres multipli\'es. Tu divises diagonalement les carr\'es
obtenus, \`a l'aide de diagonales allant du coin inf\'erieur gauche au
coin sup\'erieur droit (figure \ref{fig:ibnalbanna}).

Tu places le multiplicande au-dessus du quadrilat\`ere, en faisant 
correspondre chacune de ses positions \`a une colonne\footnote{L'écriture
du nombre s'effectue de droite à gauche (exemple : 352 s'écrira donc 253).}. 
Puis, tu places le multiplicateur \`a gauche ou \`a droite du quadrilat\`ere,
de telle sorte qu'il descende avec lui en faisant correspondre \'egalement 
chacune de ses positions \`a une ligne\footnote{L'écriture
du nombre s'effectue de bas en haut (exemple : {\tiny$\begin{array}{c}3\\5\\2\end{array}$} 
s'écrira donc {\tiny$\begin{array}{c}2\\5\\3\end{array}$}).}. Puis, tu multiplies, 
l'une apr\`es l'autre, chacune des positions du multiplicande du carr\'e 
par toutes les positions du multiplicateur, et tu poses le r\'esultat 
partiel correspondant \`a chaque position dans le carr\'e o\`u se coupent 
respectivement leur colonne et leur ligne, en pla\c{c}ant les unit\'es 
au-dessus de la diagonale et les dizaines en dessous. Puis, tu
commences \`a additionner, en partant du coin sup\'erieur gauche :
tu additionnes ce qui est entre les diagonales, sans effacer, 
en pla\c{c}ant chaque nombre dans sa position, en transf\'erant 
les dizaines de chaque somme partielle \`a la diagonale suivante et
en les ajoutant \`a ce qui y figure. 

La somme que tu obtiendras sera le r\'esultat. »}

\noindent En utilisant la m\'ethode du tableau d'Ibn al-Banna, calculer $63247\times124$ 
($= 7842628$).
\end{td}

\begin{td}\label{td:boulier}\em La division chinoise (d'après \cite{chabert}).\\
\marginpar{\footnotesize\em
\begin{fig}\label{fig:boulier}Boulier chinois.\\
\framebox[8cm]{\em
\mbox{}\hfill 8017 :\hfill
\begin{minipage}{5cm}
\includegraphics[width=5cm]{boulier.jpg}
\end{minipage}}
\end{fig}}
Dans sa version actuelle, le boulier chinois se compose d'un nombre variable de tringles
serties dans un cadre rectangulaire. Sur chacune de ces tringles, deux étages de boules
séparées par une barre transversale peuvent coulisser librement (figure \ref{fig:boulier}).
La notation des nombres repose sur le principe de la numération de position : chacune
des 2 boules du haut vaut 5 unités et chacune des 5 boules du bas vaut 1 unité. Seules
comptent les boules situées dans la région transversale.

Il existe des règles spéciales de division pour chaque diviseur de 1 à 9.
On consid\`ere ici les 7 r\`egles de division par 7 (figure \ref{fig:div7}) :

\noindent\begin{minipage}[t]{7cm}\footnotesize
\begin{enumerate}
\item {\em «~qi-yi xia jia san~»} : 7-1 ? ajouter 3 en dessous ! \\
\item {\em «~qi-er xia jia liu~»} : 7-2 ? ajouter 6 en dessous ! \\
\item {\em «~qi-san si sheng er~»} : 7-3 ? 4 reste 2 ! \\
\end{enumerate}
\end{minipage}\hfill
\begin{minipage}[t]{7cm}\footnotesize
\begin{enumerate}\setcounter{enumi}{3}
\item {\em «~qi-si wu sheng wu~»} : 7-4 ? 5 reste 5 ! \\
\item {\em «~qi-wu qi sheng yi~»} : 7-5 ? 7 reste 1 ! \\
\item {\em «~qi-liu ba sheng si~»} : 7-6 ? 8 reste 4 ! \\
\item {\em «~feng-qi jin yi~»} : 7-7 ? 1 monté ! \\
\end{enumerate}
\end{minipage}

\marginpar{\footnotesize\em
\begin{fig}\label{fig:div7}Règles de la division par 7.\\
\framebox[8cm]{\em\begin{tabular}{|c|c|c|}
\hline
Règle & Avant & Après \\
\hline
\hline
7-1   & 
\begin{tabular}{ccccc}
1 & \makebox[0.5mm]{} & 0 & 0 & 0 \\
\hline
2 & & 0 & 1 & 0
\end{tabular} &
\begin{tabular}{ccccc}
1 & \makebox[0.5mm]{} & 0 & 0 & 0 \\
\hline
2 & & 0 & 1 & 3
\end{tabular} \\
\hline
\hline
7-2   & 
\begin{tabular}{ccccc}
1 & \makebox[0.5mm]{} & 0 & 0 & 0 \\
\hline
2 & & 0 & 2 & 0
\end{tabular} &
\begin{tabular}{ccccc}
1 & \makebox[0.5mm]{} & 0 & 0 & 1 \\
\hline
2 & & 0 & 2 & 1
\end{tabular} \\
\hline
\hline
7-3   & 
\begin{tabular}{ccccc}
1 & \makebox[0.5mm]{} & 0 & 0 & 0 \\
\hline
2 & & 0 & 3 & 0
\end{tabular} &
\begin{tabular}{ccccc}
1 & \makebox[0.5mm]{} & 0 & 0 & 0 \\
\hline
2 & & 0 & 4 & 2
\end{tabular} \\
\hline
\hline
7-4   & 
\begin{tabular}{ccccc}
1 & \makebox[0.5mm]{} & 0 & 0 & 0 \\
\hline
2 & & 0 & 4 & 0
\end{tabular} &
\begin{tabular}{ccccc}
1 & \makebox[0.5mm]{} & 0 & 1 & 1 \\
\hline
2 & & 0 & 0 & 0
\end{tabular} \\
\hline
\hline
7-5   & 
\begin{tabular}{ccccc}
1 & \makebox[0.5mm]{} & 0 & 1 & 0 \\
\hline
2 & & 0 & 0 & 0
\end{tabular} &
\begin{tabular}{ccccc}
1 & \makebox[0.5mm]{} & 0 & 1 & 0 \\
\hline
2 & & 0 & 2 & 1
\end{tabular} \\
\hline
\hline
7-6   & 
\begin{tabular}{ccccc}
1 & \makebox[0.5mm]{} & 0 & 1 & 0 \\
\hline
2 & & 0 & 1 & 0
\end{tabular} &
\begin{tabular}{ccccc}
1 & \makebox[0.5mm]{} & 0 & 1 & 0 \\
\hline
2 & & 0 & 3 & 4
\end{tabular} \\
\hline
\hline
7-7   & 
\begin{tabular}{ccccc}
1 & \makebox[0.5mm]{} & 0 & 1 & 0 \\
\hline
2 & & 0 & 2 & 0
\end{tabular} &
\begin{tabular}{ccccc}
1 & \makebox[0.5mm]{} & 0 & 0 & 0 \\
\hline
2 & & 1 & 0 & 0
\end{tabular} \\
\hline
\end{tabular}}
\end{fig}}
Ces règles ne sont pas des règles logiques, 
mais de simples procédés mnémotechniques
indiquant ce qu'il convient de faire selon la situation. Leur énoncé débute par le
rappel du diviseur, ici 7, et se poursuit par l'énoncé du dividende, par exemple 3 :
7-3. Le reste de la règle indique quelles manipulations effectuées, ajouts ou retraits
de boules. Il faut également savoir que le dividende étant posé sur le boulier, on doit
appliquer les règles aux chiffres successifs du dividende, en commençant par celui dont
l'ordre est le plus élevé.
«~ajouter en dessous~» veut dire «~mettre des boules au rang 
immédiatement inférieur (à droite) au rang considéré~» et «~monter~» veut dire «~mettre des boules
au rang immédiatement supérieur (à gauche) au rang considéré~».


Pour effectuer la division d'un nombre par 7, on pose le dividende à droite sur le
boulier et le diviseur (7) à gauche. On opère sur les chiffres successifs du dividende
en commençant par celui d'ordre le plus élevé (le plus à gauche). Les règles
précédemment énoncées sont appliquées systématiquement.

Utiliser un boulier chinois pour diviser 1234 par 7 ($1234 = 176\times 7 + 2$).

\centerline{\begin{tabular}{cccccc}
1 & \makebox[1mm]{} & 0 & 0 & 0 & 0 \\
\hline
2 & & 1 & 2 & 3 & 4
\end{tabular}
$\rightarrow$
\begin{tabular}{cccccc}
1 & \makebox[1mm]{} & 0 & 1 & 1 & 0 \\
\hline
2 & & 1 & 2 & 1 & 2
\end{tabular}}

\end{td}


\begin{td}\label{td:shadok}\em Le calcul Shadok (d'après \cite{rouxel}).\\
Les cerveaux des Shadoks avaient une capacité tout à fait limitée.
Ils ne comportaient en tout que 4 cases.
Comme ils n'avaient que 4 cases, évidemment les Shadoks ne connaissaient 
pas plus de 4 mots :  {\sc ga, bu, zo et meu} (figure \ref{fig:bdShadok}).
\marginpar{\footnotesize\em 
\begin{fig}\label{fig:bdShadok}Les Shadoks : {\sc ga} {\sc bu} {\sc zo} {\sc meu}
{\em \cite{rouxel}}.\\
\noindent\framebox[8cm]{\includegraphics[width=7cm]{bdShadok.jpg}}
\end{fig}}
Etant donné qu'avec 4 mots, ils ne pouvaient pas compter plus loin que 4,
le Professeur Shadoko avait réformé tout ça :
\begin{itemize}
\item Quand il n'y a pas de Shadok, on dit {\sc ga} et on écrit {\sc ga}.
\item Quand il y a un Shadok de plus, on dit {\sc bu} et on écrit {\sc bu}.
\item Quand il y a encore un Shadok, on dit {\sc zo} et on écrit {\sc zo}.
\item Et quand il y en a encore un autre, on dit {\sc meu} et on écrit {\sc meu}.
\end{itemize}

Tout le monde applaudissait très fort et trouvait ça génial sauf le Devin Plombier 
qui disait qu'on n'avait pas idée d'inculquer à des enfants des bêtises pareilles 
et que Shadoko, il fallait le condamner. 
Il fut très applaudi aussi. Les mathématiques, 
cela les intéressait, bien sûr, mais brûler le professeur, c'était intéressant aussi, faut dire. 
Il fut décidé à l'unanimité qu'on le laisserait parler et qu'on le brûlerait après, à la récréation.
\begin{itemize}
\item Répétez avec moi : {\sc meu} {\sc zo} {\sc bu} {\sc ga}\ldots {\sc ga} {\sc bu} {\sc zo} {\sc meu}.
\item Et après! ricanait le Plombier.
\item Si je mets un Shadok en plus, évidemment, je n'ai plus assez 
de mots pour les compter, alors c'est très simple : on les jette dans une poubelle, 
et je dis que j'ai {\sc bu} poubelle. Et pour ne pas confondre avec le {\sc bu} du début, 
je dis qu'il n'y a pas de Shadok à côté de la poubelle et j'écris {\sc bu} {\sc ga}. 
{\sc bu} Shadok à côté de la poubelle: {\sc bu} {\sc bu}. Un autre : {\sc bu} {\sc zo}. Encore un autre : {\sc bu} {\sc meu}. 
\marginpar{\footnotesize\em
\begin{fig}\label{fig:calculShadok}Les 18 premiers nombres Shadok.\\
\framebox[8cm]{\begin{tabular}[t]{l@{ : }ll@{ : }ll@{ : }l}
0 & {\sc ga} 	      & 6  & {\sc bu} {\sc zo}  & 12 & {\sc meu} {\sc ga}\\
1 & {\sc bu} 	      & 7  & {\sc bu} {\sc meu} & 13 & {\sc meu} {\sc bu}\\
2 & {\sc zo} 	      & 8  & {\sc zo} {\sc ga}  & 14 & {\sc meu} {\sc zo}\\
3 & {\sc meu} 	      & 9  & {\sc zo} {\sc bu}  & 15 & {\sc meu} {\sc meu}\\
4 & {\sc bu} {\sc ga} & 10 & {\sc zo} {\sc zo}  & 16 & {\sc bu} {\sc ga} {\sc ga}\\
5 & {\sc bu} {\sc bu} & 11 & {\sc zo} {\sc meu} & 17 & {\sc bu} {\sc ga} {\sc bu}\\
\end{tabular}}
\end{fig}}
On continue. {\sc zo} poubelles et pas de Shadok à côté : {\sc zo} {\sc ga}\ldots 
{\sc meu} poubelles et {\sc meu} Shadoks à côté : {\sc meu} {\sc meu}. 
Arrivé là, si je mets un Shadok en plus, il me faut une autre poubelle. 
Mais comme je n'ai plus de mots pour compter les poubelles, je m'en débarrasse 
en les jetant dans une grande poubelle. 
J'écris {\sc bu} grande poubelle avec pas de petite poubelle 
et pas de Shadok à côté: {\sc bu} {\sc ga} {\sc ga}, et on continue\ldots {\sc bu} {\sc ga} {\sc bu}, 
{\sc bu} {\sc ga} {\sc zo}\ldots {\sc meu} {\sc meu} {\sc zo}, {\sc meu} {\sc meu} {\sc meu}.
Quand on arrive là et qu'on a trop de grandes poubelles pour pouvoir les compter, 
eh bien, on les met dans une super-poubelle, on écrit {\sc bu} {\sc ga} {\sc ga} {\sc ga}, 
et on continue\ldots (figure \ref{fig:calculShadok}).
\end{itemize}

\noindent\begin{enumerate}
\item Quels sont les entiers décimaux représentés en «~base Shadok~» 
	par les expressions suivantes ?
	\begin{enumerate}
	\item {\sc ga} {\sc ga}
	\item {\sc bu} {\sc bu} {\sc bu}
	\item {\sc zo} {\sc zo} {\sc zo} {\sc zo}
	\item {\sc meu} {\sc meu} {\sc meu} {\sc meu} {\sc meu} 
	\end{enumerate}	
\item Effectuer les calculs Shadok suivants.
	\begin{enumerate}
	\item {\sc zo} {\sc zo} {\sc meu} $+$ {\sc bu} {\sc ga} {\sc meu}
	\item {\sc meu} {\sc ga} {\sc meu} $-$ {\sc bu} {\sc meu} {\sc ga}
	\item {\sc zo} {\sc meu} {\sc meu} $\times$ {\sc bu} {\sc ga} {\sc meu}
	\item {\sc zo} {\sc zo} {\sc zo} {\sc meu} $\div$ {\sc bu} {\sc ga} {\sc zo}
	\end{enumerate}
\end{enumerate}
\end{td}

%-------------------------------------------------------------------------
\subsection{Solutions des exercices}\label{sub:solutions}
%-------------------------------------------------------------------------
\begin{description}
\item[TD \ref{td:qcm} :] QCM.\\
	Les bonnes réponses sont extraites directement
	du texte de la section \ref{contexte} : \\
	1a, 2e, 3b, 4b, 5e, 6e, 7a, 8d
\item[TD \ref{td:plage3} :] Dessins sur la plage : exécution.
	\begin{enumerate}
	\item On trace un triangle rectangle dont l'hypothénuse fait 5 pas de long
		(d'après le théorème de Pythagore : $5^2 = 3^2 + 4^2$).
	\item On a donc marché $3 + 4 + 5 = 12$ pas.
	\end{enumerate}
\item[TD \ref{td:plage4} :] Dessins sur la plage : conception.\\
	\marginpar{\footnotesize\em
	\begin{fig}\label{fig:spirale}Spirales rectangulaires.\\
	\framebox[8cm]{\includegraphics[width=7.5cm]{spirale.pdf}}
	\end{fig}}
	Imaginons l'algorithme de tracé suivant :
	\begin{enumerate}
	\item avance de 2 pas,
	\item tourne à gauche de $90^\circ$,
	\item avance de 3 pas,
	\item tourne à gauche de $90^\circ$,
	\item avance de 4 pas,
	\item tourne à gauche de $90^\circ$,
	\item avance de 5 pas,
	\item tourne à gauche de $90^\circ$,
	\item avance de 6 pas.
	\end{enumerate}
	\begin{itemize}
	\item validité : on doit au moins vérifier que la figure obtenue à toutes
		les caractéristiques recherchées : spirale rectangulaire de
		5 côtés, le plus petit côté faisant 2 pas de long et chaque 
		côté fait un pas de plus que le précédent (figure \ref{fig:spirale}).
	\item robustesse : cet algorithme suppose qu'on a suffisamment de place pour
		tracer une spirale (le dernier côté fait 6 pas de long); s'il fonctionne 
		correctement sur une plage, il ne fonctionnera certainement plus dans un
		placard.
	\item réutilisabilité : il existe une infinité de spirales rectangulaires
		qui ont les caractéristiques attendues; il suffit de penser
		à changer l'orientation initiale ou la longueur du pas par exemple. 
		On ne pourra donc pas utiliser l'algorithme tel quel dans toutes les
		configurations : il aurait fallu paramétrer l'angle de rotation et la
		longueur du pas.
	\item complexité : on peut la caractériser par le nombre de pas : 
		$2 + 3 + 4 + 5 + 6 = 20$ pas.
	\item efficacité : si la complexité se calcule en nombre de pas comme
		ci-dessus, on pourrait imaginer par exemple que la fréquence des
		pas soit plus grande («~fréquence d'horloge~») ou que 
		5 personnes prennent en charge chacune un côté de la spirale 
		pour gagner du temps («~système multi-processeurs~»).
	\end{itemize}
	\marginpar{\footnotesize\em
	\begin{rem}Loi de Moore.\\
	Gordon Earl Moore est le co-fondateur avec Robert Noyce et Andrew Grove de la société Intel en 1968 
	(fabriquant n°1 mondial de microprocesseurs).
	En 1965, il expliquait que la complexité des semiconducteurs 
	doublait tous les dix-huit mois à coût constant depuis 1959, date de leur invention. 
	En 1975, il précise sa «~première loi~» en affirmant que le nombre de transistors des microprocesseurs  
	sur une puce de silicium double tous les deux ans («~deuxième loi~»).
	
	\centerline{\includegraphics[width=8cm]{loimoore.jpg}}
	\end{rem}
	}
\item[TD \ref{td:mips} :] Puissance de calcul.\\
	L'unité de puissance est donné ici en Mips («~million d'instructions par
	seconde~»).
	\begin{enumerate}
	\item le premier micro-ordinateur de type PC : $\approx 10^{-1}$ Mips
	\item une console de jeu actuelle : $\approx 10^{4}$ Mips
	\item un micro-ordinateur actuel : $\approx 10^{4}$ Mips
	\item {\em Deeper-Blue} : l'ordinateur qui a « battu » Kasparov aux 
		échecs en 1997 : $\approx 10^{7}$ Mips
	\item le plus puissant ordinateur actuel : $\approx 10^{9}$ Mips
	\end{enumerate}
\item[TD \ref{td:stock} :] Stockage de données.
	\begin{enumerate}
	\item une page d'un livre : $\approx$ 50 lignes de 80 caractères = 4 ko
	\item une encyclopédie en 20 volumes : $\approx$ 20 volumes de 1000
	pages de 50 lignes de 80 caractères = 80 Mo (sans images !)
	\item une photo couleur : $\approx$ de 1Mo à 100 Mo selon la qualité
	\item une heure de vidéo : $\approx$ de 500 Mo à 2 Go 
		selon la qualité vidéo (DivX, MPEG-2\ldots)
	\item une minute de son : $\approx$ de 1 Mo (MP3) à 10 Mo 
		selon la qualité du son
	\item une heure de son : $\approx$ de 60 Mo à 600 Mo 
		selon la qualité du son
	\end{enumerate}
\newpage
\item[TD \ref{td:tortue} :] Tracés de polygones réguliers.
	\begin{enumerate}
	\item Pentagone régulier de 10 pas de côté.\\
		\begin{minipage}{6cm}
		\begin{enumerate}
		\item avance de 10 pas,
		\item tourne à gauche de $(360/5)^\circ$,
		\item avance de 10 pas,
		\item tourne à gauche de $(360/5)^\circ$,
		\item avance de 10 pas,
		\item tourne à gauche de $(360/5)^\circ$,
		\item avance de 10 pas,
		\item tourne à gauche de $(360/5)^\circ$,
		\item avance de 10 pas,
		\item tourne à gauche de $(360/5)^\circ$.
		\end{enumerate}
		\end{minipage}
		\hfill
		\begin{minipage}{7cm}
		On remarque qu'on a effectué 5 fois de suite les 2 instructions 
		suivantes :
		\begin{enumerate}
		\item avance de 10 pas,
		\item tourne à gauche de $(360/5)^\circ$.
		\end{enumerate}
		Pour simplifier, on écrira plutôt :\\
		Répète 5 fois de suite les 2 instructions
		\begin{enumerate}
		\item avance de 10 pas,
		\item tourne à gauche de $(360/5)^\circ$.
		\end{enumerate}
		C'est ce que nous ferons dans les exemples suivants.
		\end{minipage}
	\item Hexagone régulier de 10 pas de côté.\\
		Répète 6 fois de suite les 2 instructions 
		\begin{enumerate}
		\item avance de 10 pas,
		\item tourne à gauche de $(360/6)^\circ$,
		\end{enumerate}
	\item Octogone régulier de 10 pas de côté.\\
		Répète 8 fois de suite les 2 instructions 
		\begin{enumerate}
		\item avance de 10 pas,
		\item tourne à gauche de $(360/8)^\circ$,
		\end{enumerate}
	\item Polygone régulier de $n$ côtés de 10 pas chacun.\\
		Répète $n$ fois de suite les 2 instructions 
		\begin{enumerate}
		\item avance de 10 pas,
		\item tourne à gauche de $(360/n)^\circ$,
		\end{enumerate}
	\end{enumerate}
\newpage
\item[TD \ref{td:russe} :] Multiplication «~à la russe~».
	\begin{enumerate}
	\item \begin{tabular}[t]{|r|r|r|r|}
		\hline
		multiplicande & multiplicateur & reste      & somme partielle\\
		$M \times 2$  & $m \div 2$     & $m \bmod 2$ &       \\
		\hline
		 96  & 64 &   0   &   $(0\times 96) + 0$ \\
		192  & 32 &   0   &   $(0\times 192) + 0$ \\
		384  & 16 &   0   &   $(1\times 384) + 0$ \\
		768  &  8 &   0   &   $(0\times 768) + 0$  \\
	       1536  &  4 &   0   &   $(0\times 1536) + 0$ \\
	       3072  &  2 &   0   &   $(0\times 3072) + 0$ \\
	       6144  &  1 &   1   &   $(1\times 6144) + 0$ \\
		\hline
		\multicolumn{3}{|r|}{$64 \times 96 =$} & 6144\\
		\hline
		\end{tabular}

	\item \begin{tabular}[t]{|r|r|r|r|}
		\hline
		multiplicande & multiplicateur & reste      & somme partielle\\
		$M \times 2$  & $m \div 2$     & $m \bmod 2$ &       \\
		\hline
		239  & 45 &   1   &   $(1\times 239) + 0$ \\
		478  & 22 &   0   &   $(0\times 478) + 239$ \\
		956  & 11 &   1   &   $(1\times 956) + 239$ \\
	       1912  &  5 &   1   &   $(1\times 1912) + 1195$  \\
	       3824  &  2 &   0   &   $(0\times 3824) + 3107$ \\
	       7648  &  1 &   1   &   $(1\times 7648) + 3107$ \\
		\hline
		\multicolumn{3}{|r|}{$45 \times 239 =$} & 10755\\
		\hline
		\end{tabular}
	\end{enumerate}
\item[TD \ref{td:ibnalbanna} :] Multiplication arabe : $63247\times124 = 7842628$.\\
	\mbox{}\hspace*{0.9cm}{\includegraphics[width=7cm]{ibnalbanna1.pdf}} 
\item[TD \ref{td:boulier} :] Division chinoise : $1234 \div 7$ 
	($1234 = 176\times 7 +2$).
	\marginpar{\footnotesize\em
	\begin{rem}Boulier chinois : division par 3.
	\begin{enumerate}
	\item {\em «~san-yi sanshi-yi~»} : trois-un ? trente et un ! \\
		(on pose 3 \`a la place du 1, et on ajoute 1 \`a droite)
	\item {\em «~san-er liushi-er~»} : trois-deux ? soixante deux ! \\
		(on pose 6 \`a la place du 2, et on ajoute 2 \`a droite)
	\item {\em «~feng san jin yi-shi~»} : trois-trois ? dizaine mont\'ee ! \\
		(on pose 0 \`a la place du 3, et on ajoute 1 \`a gauche)
	\end{enumerate}
	Effectuer la division $2271 \div 3$ avec le boulier chinois.
	\end{rem}}
	
	\noindent\begin{tabular}{|cccccc|}
	\hline
	  &                 & $\downarrow$ & & & \\
	1 & \makebox[1mm]{} & 0 & 0 & 0 & 0 \\
	\hline
	2 & & 1 & 2 & 3 & 4 \\
	\hline
	\end{tabular}
	$\rightarrow$ 7-1 $\rightarrow$ 
	\begin{tabular}{|cccccc|}
	\hline
	  &                 & & $\downarrow$ & & \\
	1 & \makebox[1mm]{} & 0 & 1 & 0 & 0 \\
	\hline
	2 & & 1 & 0 & 3 & 4\\
	\hline
	\end{tabular}
	$\rightarrow$ 7-5 $\rightarrow$ 
	\begin{tabular}{|cccccc|}
	\hline
	  &                 & & & $\downarrow$ & \\
	1 & \makebox[1mm]{} & 0 & 1 & 0 & 0 \\
	\hline
	2 & & 1 & 2 & 4 & 4\\
	\hline
	\end{tabular}
	
	\noindent\begin{tabular}{|cccccc|}
	\hline
	  &                 & & & $\downarrow$ & \\
	1 & \makebox[1mm]{} & 0 & 1 & 0 & 0 \\
	\hline
	2 & & 1 & 2 & 4 & 4\\
	\hline
	\end{tabular}
	$\rightarrow$ 7-4 $\rightarrow$ 
	\begin{tabular}{|cccccc|}
	\hline
	  &                 & & & & $\downarrow$ \\
	1 & \makebox[1mm]{} & 0 & 1 & 1 & 1 \\
	\hline
	2 & & 1 & 2 & 0 & 4\\
	\hline
	\end{tabular}
	$\rightarrow$ 7-7 $\rightarrow$ 
	\begin{tabular}{|cccccc|}
	\hline
	  &                 & & & & \\
	1 & \makebox[1mm]{} & 0 & 1 & 1 & 0 \\
	\hline
	2 & & 1 & 2 & 1 & 2\\
	\hline
	\end{tabular}
\item[TD \ref{td:shadok} :] Calcul Shadok.\\
	Le système Shadok est un système de numération en base 4 :\\
	{\sc ga} = 0, {\sc bu} = 1, {\sc zo} = 2 et {\sc meu} = 3.
	\marginpar{\footnotesize\em
	\begin{rem}
	Un entier positif en base $b$ est représenté par une suite de
	chiffres {$(r_nr_{n-1}\ldots r_1r_0)_b$}
	où les $r_i$ sont des chiffres de la base $b$ ($0\leq r_i < b$).
	Ce nombre a pour valeur:
	$${r_nb^n + r_{n-1}b^{n-1} + \ldots + r_1b^1 + r_0b^0 = 
        \sum^{i=n}_{i = 0} r_ib^i}$$
	
	Un nombre fractionnaire (nombre avec des chiffres après la virgule :
	{$(r_nr_{n-1}\ldots r_1r_0.r_{-1}r_{-2}\ldots)_b$})
	est défini sur un sous-en\-semble borné, incomplet et fini des rationnels.
	Un tel nombre a pour valeur :
	$${r_nb^n + r_{n-1}b^{n-1} + \ldots + r_0b^0 + r_{-1}b^{-1} + r_{-2}b^{-2} + \ldots}$$
	En pratique, le nombre de chiffres après la virgule est limité par la taille physique
	en machine.
	$${(r_nr_{n-1}\ldots r_1r_0.r_{-1}r_{-2}\ldots r_{-m})_b = \sum_{i=-m}^{i=n} r_ib^i}$$
	\end{rem}}
	\begin{enumerate}
	\item   \begin{enumerate}
		\item {\sc ga} {\sc ga} 
			$= (00)_4 = 0$
		\item {\sc bu} {\sc bu} {\sc bu} 
			$= (111)_4 = 1\cdot 4^2 + 1\cdot 4^1 + 1\cdot 4^0 = 16 + 4 + 1 = 21$
		\item {\sc zo} {\sc zo} {\sc zo} {\sc zo} 
			$= (2222)_4 = 2\cdot 4^3 + 2\cdot 4^2 + 2\cdot 4^1 + 2\cdot 4^0 = 128 + 32 + 8 + 2 = 170$
		\item {\sc meu} {\sc meu} {\sc meu} {\sc meu} {\sc meu} 
			$= (33333)_4 = 3 \cdot 4^4 + 3\cdot 4^3 + 3\cdot 4^2 + 3\cdot 4^1 + 3\cdot 4^0 = 768 + 192 + 48 + 12 + 3 =
			1023$
		\end{enumerate}	
	\item	\begin{enumerate}
		\item {\sc zo} {\sc zo} {\sc meu} $+$ {\sc bu} {\sc ga} {\sc meu} 
			$= (223)_4 + (103)_4 = (332)_4 = 43 + 19 = 62$
		\item {\sc meu} {\sc ga} {\sc meu} $-$ {\sc bu} {\sc meu} {\sc ga} 
			$= (303)_4 - (130)_4 = (113)_4 = 51 - 28 = 23$
		\item {\sc zo} {\sc meu} {\sc meu} $\times$ {\sc bu} {\sc ga} {\sc meu} 
			$= (233)_4 \times (103)_4 = (31331)_4 = 47 \times 19 = 893$
		\item {\sc zo} {\sc zo} {\sc zo} {\sc meu} $\div$ {\sc bu} {\sc ga} {\sc zo} 
			$= (2223)_4 \div (102)_4 = (21)_4 = 171 \div 18 = 9$
		\end{enumerate}
	\end{enumerate}
\end{description}

%-------------------------------------------------------------------------
\newpage
\section{Annexes}
%-------------------------------------------------------------------------

%-------------------------------------------------------------------------
\subsection{Lettre de Jacques Perret}\label{annexe:perret}
%-------------------------------------------------------------------------
\marginpar{\footnotesize\em
\begin{fig}\label{fig:eniac}Le premier ordinateur (1946) : ENIAC (Electronic Numerical Integrator 
Analyser and Computer).\\
\framebox[8cm]{\includegraphics[width=6cm]{eniac2.jpg}}
\end{fig}
\begin{fig}\label{fig:micros}Premiers micro-ordinateurs.\\
\framebox[8cm]{\begin{minipage}{7.5cm}
\includegraphics[width=3.25cm]{apple2.jpg}\hfill
\includegraphics[width=3.25cm]{zx81.jpg}

\mbox{}\hfill Apple II (1977) \hfill ZX 81 (1981) \hfill \mbox{}

\includegraphics[width=3.25cm]{ibmpc.jpg}\hfill
\includegraphics[width=3.25cm]{macintosh.jpg}

\mbox{}\hfill IBM PC (1981)\hfill Macintosh (1984)\hfill \mbox{}
\end{minipage}}
\end{fig}}
Au printemps de 1955, IBM France s'apprêtait à construire dans ses ateliers 
de Corbeil-Essonnes (consacrés jusque-là au montage des machines mécanographiques 
--- tabulatrices, trieuses, \ldots --- de technologie électromécanique) 
les premières machines électroniques destinées au traitement de l'information. 
Aux États-Unis ces nouvelles machines étaient désignées sous le vocable 
Electronic Data Processing System. Le mot «~computer~» était plutôt réservé 
aux machines scientifiques et se traduisait aisément en «~calculateur~» ou «~calculatrice~».
Sollicité par la direction de l'usine de Corbeil-Essonnes, François Girard, 
alors responsable du service promotion générale publicité, 
décida de consulter un de ses anciens maîtres, Jacques Perret, professeur de philologie 
latine à la Sorbonne. A cet effet il écrit une lettre à la signature de C. de Waldner, 
président d'IBM France. Il décrit sommairement la nature et les fonctions des 
nouvelles machines. Il accompagne sa lettre de brochures illustrant les machines 
mécanographiques. Le 16 avril, le professeur Perret lui répond.
L'ordinateur IBM 650 peut commencer sa carrière.
Protégé pendant quelques mois par IBM France, le mot fut rapidement adopté
par un public de spécialistes, de chefs d'entreprises et par l'administration. 
IBM décida de le laisser dans le domaine public
(d'après le site de la $10^{\grave eme}$ semaine de la langue française et de la francophonie
{\tt www.semainelf.culture.fr/site2005/dixmots/}).

{\em
Le 16 IV 1955

Cher Monsieur,

Que diriez-vous d'«~ordinateur~» ? C'est un mot correctement formé, qui se trouve même dans le
Littré comme adjectif désignant Dieu qui met de l'ordre dans le monde. Un mot de ce genre a
l'avantage de donner aisément un verbe «~ordiner~», un nom d'action «~ordination~». L'inconvénient
est que «~ordination~» désigne une cérémonie religieuse ; mais les deux champs de signification
(religion et comptabilité) sont si éloignés et la cérémonie d'ordination connue, je crois, de si peu de
personnes que l'inconvénient est peut-être mineur. D'ailleurs votre machine serait «~ordinateur~» (et
non ordination) et ce mot est tout à fait sorti de l'usage théologique. «~Systémateur~» serait un
néologisme, mais qui ne me paraît pas offensant ; il permet «~systématisé~» ; - mais «~système~» ne
me semble guère utilisable - «~combinateur~» a l'inconvénient du sens péjoratif de «~combine~» ;
«~combiner~» est usuel donc peu capable de devenir technique ; «~combination~» ne me paraît guère
viable à cause de la proximité de «~combinaison~». Mais les Allemands ont bien leurs «~combinats~»
(sorte de trusts, je crois), si bien que le mot aurait peut-être des possibilités autres que celles
qu'évoque «~combine~».

\marginpar{\footnotesize\em
\begin{fig}\label{fig:intel86}Du 8086 (1978) au Core 2 (2006) d'Intel.\\
\framebox[8cm]{\begin{minipage}{7.5cm}
\includegraphics[width=3.25cm]{8086.jpg}\hfill
\includegraphics[width=3.25cm]{80486.jpg}

\mbox{}\hfill 8086 \hfill 80486 \hfill \mbox{}

\includegraphics[width=3.25cm]{pentium4.jpg}\hfill
\includegraphics[width=3.25cm]{core2.jpg}

\mbox{}\hfill Pentium 4 \hfill Core Duo \hfill \mbox{}
\end{minipage}}
\end{fig}

\begin{fig}\label{fig:derniersnes}Micro-ordinateurs récents.\\
\framebox[8cm]{\begin{minipage}{7.5cm}
\includegraphics[width=3.25cm]{pda.jpg}\hfill
\includegraphics[width=3.25cm]{tablette.jpg}

\mbox{}\hfill PDA \hfill Tablette \hfill \mbox{}

\includegraphics[width=2.5cm]{iphone.jpg}\hfill
\includegraphics[width=3.25cm]{wii.jpg}

\mbox{}\hfill iPhone \hfill Wii \hfill \mbox{}
\end{minipage}}
\end{fig}
}
«~Congesteur~», «~digesteur~» évoquent trop «~congestion~» et «~digestion~».
«~Synthé\-ti\-seur~» ne me
paraît pas un mot assez neuf pour designer un objet spécifique, déterminé comme votre machine.
En relisant les brochures que vous m'avez données, je vois que plusieurs de vos appareils sont
désignés par des noms d'agent féminins (trieuse, tabulatrice). «~Ordinatrice~» serait parfaitement
possible et aurait même l'avantage de séparer plus encore votre machine du vocabulaire de la
théologie.
Il y a possibilité aussi d'ajouter à un nom d'agent un complément : «~ordinatrice d'éléments
complexes~» ou un élément de composition, par ex. : «~sélecto-sys\-té\-ma\-teur~». - «~Sélecto-
ordinateur~» a l'inconvénient de 2 «~o~» en hiatus, comme «~électro-ordinatrice~».
Il me semble que je pencherais pour «~ordinatrice électroni\-que~». Je souhaite que ces suggestions
stimulent, orientent vos propres facultés d'invention. N'hésitez pas à me donner un coup de
téléphone si vous avez une idée qui vous paraisse requérir l'avis d'un philologue.

Vôtre.

J. Perret
}


%-------------------------------------------------------------------------
\newpage
\subsection{Exemple de questionnaire d'évaluation}\label{annexe:eval}
%-------------------------------------------------------------------------
$$\begin{tabular}{|p{11cm}|ccccc|}
\hline
\bf Proposition & \bf 1 & \bf 2 & \bf 3 & \bf 4 & \bf $\times$ \\
\hline
\hline
Vos connaissances antérieures étaient suffisantes pour suivre ce cours & $\Box$ & $\Box$ & $\Box$ & $\Box$ & $\Box$ \\
\hline
Les objectifs du cours ont été clairement définis & $\Box$ & $\Box$ & $\Box$ & $\Box$ & $\Box$ \\
\hline
Vous êtes satisfait des supports de cours fournis & $\Box$ & $\Box$ & $\Box$ & $\Box$ & $\Box$ \\
\hline
Vous êtes satisfait de l'équilibre cours/TD/TP & $\Box$ & $\Box$ & $\Box$ & $\Box$ & $\Box$ \\
\hline
Vous êtes satisfait de la présentation de ce cours (clarté d'expression\ldots) & $\Box$ & $\Box$ & $\Box$ & $\Box$ & $\Box$ \\
\hline
Si vous avez plusieurs professeurs, la cohérence de l'enseignement vous a 
	semblé assurée & $\Box$ & $\Box$ & $\Box$ & $\Box$ & $\Box$ \\
\hline
Le rythme de travail vous a permis de suivre et de comprendre & $\Box$ & $\Box$ & $\Box$ & $\Box$ & $\Box$ \\
\hline
Le temps alloué à cette matière vous a semblé satisfaisant & $\Box$ & $\Box$ & $\Box$ & $\Box$ & $\Box$ \\
\hline
Cette matière a nécessité beaucoup de travail personnel & $\Box$ & $\Box$ & $\Box$ & $\Box$ & $\Box$ \\
\hline
Vous avez eu l'impression de progresser & $\Box$ & $\Box$ & $\Box$ & $\Box$ & $\Box$ \\
\hline
Les contrôles sont adaptés à l'objectif et au niveau du cours & $\Box$ & $\Box$ & $\Box$ & $\Box$ & $\Box$ \\
\hline
Vous êtes satisfait du mode d'évaluation & $\Box$ & $\Box$ & $\Box$ & $\Box$ & $\Box$ \\
\hline
Après les contrôles, les enseignants fournissent des commentaires qui aident
	à mieux maitriser la matière & $\Box$ & $\Box$ & $\Box$ & $\Box$ & $\Box$ \\
\hline
Vous êtes satisfait des conditions de travail 
	(salles de cours, matériel utilisé\ldots) & $\Box$ & $\Box$ & $\Box$ & $\Box$ & $\Box$ \\
\hline
Respect du programme pédagogique & $\Box$ & $\Box$ & $\Box$ & $\Box$ & $\Box$ \\
\hline
\hline
\multicolumn{6}{|c|}{{\bf 1:}Très satisfait , {\bf 2:}Satisfait , {\bf 3:}Peu satisfait , 
                   {\bf 4:}Pas satisfait ; {\bf $\times$:}Ne me concerne pas}\\
\hline
\end{tabular}$$

%-------------------------------------------------------------------------
\newpage
\subsection{Exemple de planning}\label{annexe:planning}
%-------------------------------------------------------------------------
Les enseignements d'Informatique S1 s'étalent sur 15 semaines :
de la semaine 2 à la semaine 16 dans le planning ci-dessous.
$$\begin{tabular}{|l|p{4cm}|p{4cm}|p{4cm}|}
\hline
\bf Semaine & \bf Cours & \bf TD & \bf Thème contrôle \\
\hline
\bf 1 & --- & --- & --- \\
\hline
\bf 2 & introduction générale\newline instructions de base & ---  & --- \\
\hline
\bf 3 & --- & environnement de pro\-grammation\newline instructions de base & ---  \\
\hline
\bf 4 & {\bf contrôle d'autoformation}\newline instructions de base & --- & calculs booléens\\
\hline
\bf 5 & --- & instructions de base & --- \\
\hline
\bf 6 & instructions de base \newline {\bf DS1} & --- & \mbox{}\newline instructions de base \\
\hline
\bf 7 & --- & instructions de base & --- \\
\hline
\bf 8 & procédures et fonctions & --- & --- \\
\hline
\bf 9 & --- & procédures et fonctions & --- \\
\hline
\bf 10 & {\bf contrôle d'autoformation}\newline procédures et fonctions & --- & codage des
nombres\\
\hline
\bf 11 & --- & procédures et fonctions & --- \\
\hline
\bf 12 & procédures et fonctions & --- & --- \\
\hline
\bf 13 & --- & procédures et fonctions & --- \\
\hline
\bf 14 & {\bf DS2}\newline structures linéaires & --- & procédures et fonctions\\
\hline
\bf 15 & --- &  structures linéaires & ---\\
\hline
\bf 16 & structures linéaires\newline {\bf contrôle d'autoformation} & --- & \mbox{}\newline recherche d'un élément \\
\hline
\end{tabular}$$

%-------------------------------------------------------------------------
\newpage
\subsection{Informatique à l'ENIB}\label{annexe:enib}
%-------------------------------------------------------------------------
\marginpar{\footnotesize ENIB: {\tt www.enib.fr}}
\marginpar{\includegraphics[width=8cm]{enibHaut-1280.jpg}}

%-------------------------------------------------------------------------
\subsection*{Enseignements de premier cycle}
%-------------------------------------------------------------------------
Les 4  semestres du premier cycle (S1 à S4) sont communs à
tous les élèves de l'ENIB. Les enseignements d'informatique
sont précisés dans le tableau ci-dessous.

$$\begin{tabular}{llr}
\hline
Semestre & Thème & Horaires \\
\hline
S1 & Algorithmique                 & 45 h \\
S2 & Méthode de développement      & 45 h \\
\hline
S3 & Programmation procédurale     & 45 h \\
S3 & Programmation par objets      & 45 h \\
S4 & Programmation pour l'embarqué & 45 h \\
\hline
\end{tabular}$$


%-------------------------------------------------------------------------
\subsection*{Enseignements du cycle ingénieur}
%-------------------------------------------------------------------------
Les 2 premiers semestres du cycle ingénieur (S5 et S6) sont communs à
tous les élèves de l'ENIB.

$$\begin{tabular}{llr}
\hline
Semestre & Thème & Horaires \\
\hline
S5  & Programmation par objets               & 45 h \\
S6  & Programmation par objets               & 45 h \\
S6  & Modèles pour l'ingénierie des systèmes & 45 h \\
S6  & Bases de données                       & 22 h\\
\hline
\end{tabular}$$

Au cours du $3^{\grave eme}$ semestre du cycle ingénieur (S7), 
les élèves choisissent une option parmi trois : électronique, informatique 
ou mécatronique. Les modules de l'option informatique
sont listés dans le tableau ci-dessous. Le $4^{\grave eme}$ semestre (S8) 
est consacré à un Projet Professionnalisant en Equipe (PPE).

$$\begin{tabular}{llr}
\hline
Semestre & Thème & Horaires \\
\hline
S7  & Systèmes embarqués 1                 & 45 h \\
S7  & Systèmes embarqués 2                 & 45 h \\
S7  & Réseaux                              & 45 h \\
S7  & Administration Systèmes et Réseaux   & 45 h \\
S7  & Génie Logiciel                       & 45 h \\
S7  & Systèmes d'Information               & 45 h \\
S7  & Interfaces Homme-Machine             & 45 h \\
S7  & Applications Réparties               & 45 h \\
\hline
S8  & Projet Professionnalisant en Equipe  & 350 h \\
\hline
\end{tabular}$$

Au cours du $5^{\grave eme}$ semestre du cycle ingénieur (S9),
les étudiants doivent suivre 5 modules scientifiques,
3 modules de type SHS et 1 module d'ouverture. 
les modules à caractère informatique sont donnés dans le tableau ci-dessous.
Le dernier semestre 
est consacré au stage en entreprise. 

$$\begin{tabular}{llr}
\hline
Semestre & Thème & Horaires \\
\hline
S9  & Systèmes distribués                       & 45 h \\
S9  & Contrôle adaptatif                        & 45 h \\
S9  & Styles de programmation                   & 45 h \\
S9  & Simulation                                & 45 h \\
S9  & Réalité Virtuelle                         & 45 h \\
\hline
S10 & Stage en entreprise & 20 sem \\
\hline
\end{tabular}$$


%-------------------------------------------------------------------------
\subsection*{Recherches en informatique}
%-------------------------------------------------------------------------
L'ENIB accueille trois laboratoires de recherche correspondant aux trois options :
électronique, informatique et mécatronique.

Le laboratoire de recherche en informatique de l'ENIB est le {\sc LISyC} 
(Laboratoire d'Informatique des Systèmes Complexes).
\marginpar{{\sc LISyC} : \footnotesize\tt www.lisyc.univ-brest.fr}
Le {\sc LISyC} est à Brest une Equipe d'Accueil (EA 3883) commune à l'Université de Bretagne Occidentale (UBO), 
à l'Ecole Nationale d'Ingénieurs de Brest (ENIB) et à l'Ecole Nationale Supérieure des Ingénieurs 
des Etudes des Techniques d'Armement (ENSIETA).

Le {\sc LISyC} a fait sienne l'idée selon laquelle 
{\em
comprendre et maîtriser les comportements des systèmes complexes,
naturels ou artificiels, constituent pour les scientifiques le défi majeur
du $21^{\grave eme}$ siècle.
}
Il compte actuellement 40 enseignants-chercheurs permanents et 20 doctorants
regroupés en 4 équipes STIC et 1 équipe SHS qui explorent de manière complémentaire
cette problématique :
$$\begin{tabular}{l@{ : }l}
{\sc ARéVi} & Ateliers de Réalité Virtuelle \\
{\sc EBV} & Ecosystémique et Biologie Virtuelles \\
{\sc IdM} & Ingénierie des Modèles \\
{\sc SARA} & Simulation, Apprentissages, Représentations, Action \\
{\sc SuSy} & Sûreté des Systèmes
\end{tabular}$$
Les équipes {\sc ARéVi}, {\sc EBV} et {\sc SARA} sont localisées au {\sc CERV} :
le Centre Européen de Réalité Virtuelle, établissement de l'ENIB ouvert en juin 2004.
\marginpar{\footnotesize\tt {\sc CERV} : www.cerv.fr}
\marginpar{\includegraphics[width=8cm]{cervFacade-1280.jpg}}


%-------------------------------------------------------------------------
\newpage
\section*{Liste des exercices}\label{sub:listexo}
\addcontentsline{toc}{section}{Liste des exercices}
%-------------------------------------------------------------------------
\subsubsection*{Exercices de TD}
\noindent
\exercice{td:plage1}{Dessins sur la plage : exécution (1).}\\
\exercice{td:plage2}{Dessins sur la plage : conception (1).}\\
\exercice{td:prop}{Propriétés d'un algorithme.}\\
\exercice{td:octets}{Unités d'information.}\\
\exercice{td:python}{Première utilisation de {\sc Python}.}\\
\exercice{td:erreur}{Erreur de syntaxe en {\sc Python}.}\\
\exercice{td:plage5}{Dessins sur la plage : persévérance.}\\
\exercice{td:autonomie}{Autonomie.}\\
\exercice{td:site}{Site {\sc Web} d'Informatique S1.}\\
\exercice{td:attention1}{Exemple de contrôle d'attention (1).}\\
\exercice{td:TD}{Exemple de contrôle de TD.}\\
\exercice{td:bool}{Exemple de contrôle d'autoformation (1).}\\
\exercice{td:competences}{Exemple de contrôle des compétences.}\\
\exercice{td:controles}{Nombre de contrôles.}\\
\exercice{td:negation}{Exemple de contrôle d'autoformation (2).}\\
\exercice{td:attention2}{Exemple de contrôle d'attention (2).}\\
\exercice{td:exercices}{Nombres d'exercices de TD.}\\
\exercice{td:labo}{Environnement de travail.}

\subsubsection*{Exercices complémentaires}
\noindent
\exercice{td:qcm}{QCM.}\\
\exercice{td:plage3}{Dessins sur la plage : exécution (2).}\\
\exercice{td:plage4}{Dessins sur la plage : conception (2).}\\
\exercice{td:mips}{Puissance de calcul.}\\
\exercice{td:stock}{Stockage de données.}\\
\exercice{td:tortue}{Tracés de polygones réguliers.}\\
\exercice{td:russe}{La multiplication «~à la russe~».}\\
\exercice{td:ibnalbanna}{La multiplication arabe.}\\
\exercice{td:boulier}{La division chinoise.}\\
\exercice{td:shadok}{Le calcul Shadok.}
%-------------------------------------------------------------------------

